\section{Homogeneous spectrum of a sheaf of graded algebras}
\label{section:II.3}

\subsection{Homogeneous spectrum of a quasi-coherent graded
$\mathcal{O}_Y$-algebra.}
\label{subsection:II.3.1}
\label{II.3.1}

\begin{env}[3.1.1]
\label{II.3.1.1}
Let $Y$ be a prescheme, $\mathcal{S}$ a graded $\mathcal{O}_Y$-algebra, and
$\mathcal{M}$ a graded $\mathcal{S}$-module. If $\mathcal{S}$ is
\emph{quasi-coherent}, each of its homogeneous components $\mathcal{S}_n$ is
a \emph{quasi-coherent} $\mathcal{O}_Y$-module, being the image of
$\mathcal{S}$ under a homomorphism of $\mathcal{S}$ into itself
((\hyperref[I.1.3.8]{I, 1.3.8}) and
(\hyperref[I.1.3.9]{I, 1.3.9})); likewise, if $\mathcal{M}$ is
quasi-coherent as an $\mathcal{O}_Y$-module, so are its homogeneous
components $\mathcal{M}_n$, and conversely. If $d$ is an integer $>0$, we
denote by $\mathcal{S}^{(d)}$ the direct sum of the
$\mathcal{O}_Y$-modules $\mathcal{S}_{nd}$ ($n\in\mathbf{Z}$), which is
quasi-coherent if $\mathcal{S}$ is
(\hyperref[I.1.3.9]{I, 1.3.9}); for every integer $k$ such that
$0\leq k\leq d-1$, we denote by $\mathcal{M}^{(d,k)}$ (or
$\mathcal{M}^{(d)}$ for $k=0$) the direct sum of the
$\mathcal{M}_{nd+k}$ ($n\in\mathbf{Z}$), which is a graded
$\mathcal{S}^{(d)}$-module, quasi-coherent when $\mathcal{S}$ and
$\mathcal{M}$ are quasi-coherent (\hyperref[I.9.6.1]{I, 9.6.1}). We
denote by $\mathcal{M}(n)$ the graded $\mathcal{S}$-module such that
$(\mathcal{M}(n))_k=\mathcal{M}_{n+k}$ for every $k\in\mathbf{Z}$; if
$\mathcal{S}$ and $\mathcal{M}$ are quasi-coherent, $\mathcal{M}(n)$ is a
quasi-coherent graded $\mathcal{S}$-module
(\hyperref[I.9.6.1]{I, 9.6.1}).

We will say that $\mathcal{M}$ is a graded $\mathcal{S}$-module
\emph{of finite type} (resp. that it admits a \emph{finite presentation}) if,
for every $y\in Y$, there exist an open neighbourhood $U$ of $y$ and integers
$n_i$ (resp. integers $m_i$ and $n_j$) such that there is a surjective
homomorphism of degree~$0$
\[
  \bigoplus_{i=1}^r(\mathcal{S}(n_i)|U)\longrightarrow\mathcal{M}|U
\]
(resp. such that $\mathcal{M}|U$ is isomorphic to the cokernel of a
homomorphism of degree~$0$
\[
  \bigoplus_{i=1}^r(\mathcal{S}(m_i)|U)
  \longrightarrow\bigoplus_{j=1}^s(\mathcal{S}(n_j)|U)).
\]

Let $U$ be an affine open subset of $Y$ and $A=\Gamma(U,\mathcal{O}_Y)$ its
ring; by hypothesis, the graded $(\mathcal{O}_Y|U)$-algebra $\mathcal{S}|U$
is isomorphic to $\widetilde{S}$, where $S=\Gamma(U,\mathcal{S})$ is a graded
$A$-algebra

\oldpage[II]{50}

(\hyperref[I.1.4.3]{I, 1.4.3}); set
$X_U=\operatorname{Proj}(\Gamma(U,\mathcal{S}))$. Let $U'\subset U$ be a
second affine open subset of $Y$, $A'$ its ring, and $j:U'\to U$ the
canonical injection, which corresponds to the restriction homomorphism
$A\to A'$; we have $\mathcal{S}|U'=j^*(\mathcal{S}|U)$, and consequently
$S'=\Gamma(U',\mathcal{S})$ is canonically identified with $S\otimes_A A'$
(\hyperref[I.1.6.4]{I, 1.6.4}). We conclude from this
(\hyperref[II.2.8.10]{2.8.10}) that $X_{U'}$ is canonically identified
with $X_U\times_U U'$, and consequently also with $f_U^{-1}(U')$, where $f_U$
denotes the structural morphism $X_U\to U$
(\hyperref[I.4.4.1]{I, 4.4.1}). We denote by $\sigma_{U',U}$ the
canonical isomorphism $f_U^{-1}(U')\xrightarrow{\sim}X_{U'}$ thus defined,
and by $\rho_{U',U}$ the open immersion $X_{U'}\to X_U$ obtained by composing
$\sigma_{U',U}^{-1}$ and the canonical injection $f_U^{-1}(U')\to X_U$. It is
immediate that if $U''\subset U'$ is a third affine open subset of $Y$, we
have $\rho_{U'',U}=\rho_{U',U}\circ\rho_{U'',U'}$.
\end{env}

\begin{proposition}[3.1.2]
\label{II.3.1.2}
Let $Y$ be a prescheme. For every quasi-coherent $\mathcal{O}_Y$-algebra
$\mathcal{S}$ graded in positive degrees, there exists a prescheme $X$ over
$Y$, unique up to $Y$-isomorphism, having the following property: if $f$ is
the structural morphism $X\to Y$, then for every affine open subset $U$ of
$Y$ there exists an \emph{isomorphism} $\eta_U$ of the induced prescheme
$f^{-1}(U)$ onto $X_U=\operatorname{Proj}(\Gamma(U,\mathcal{S}))$ such that,
if $V$ is a second affine open subset of $Y$ contained in $U$, the diagram
\[
\label{II.3.1.2.1}
  \xymatrix{
    f^{-1}(V)\ar[r]^{\eta_V}\ar[d]
    & X_V\ar[d]^{\rho_{V,U}}\\
    f^{-1}(U)\ar[r]_{\eta_U}
    & X_U
  }
\tag{3.1.2.1}
\]
is commutative.
\end{proposition}

\begin{proof}
For two affine open subsets $U$, $V$ of $Y$, let $X_{U,V}$ be the prescheme
induced by $X_U$ on $f_U^{-1}(U\cap V)$; we are going to define a
$Y$-isomorphism $\theta_{U,V}:X_{V,U}\xrightarrow{\sim}X_{U,V}$. To this end,
consider an affine open subset $W\subset U\cap V$: by composing the
isomorphisms
\[
  f_U^{-1}(W)\xrightarrow{\sigma_{W,U}}X_W
  \xrightarrow{\sigma_{W,V}^{-1}}f_V^{-1}(W),
\]
we obtain an isomorphism $\tau_W$, and we verify immediately that if
$W'\subset W$ is an affine open subset, $\tau_{W'}$ is the restriction of
$\tau_W$ to $f_U^{-1}(W')$; the $\tau_W$ are therefore indeed the
restrictions of a $Y$-isomorphism $\theta_{V,U}$. Moreover, if $U$, $V$, $W$
are three affine open subsets of $Y$, and $\theta'_{U,V}$, $\theta'_{V,W}$
and $\theta'_{U,W}$ are the restrictions of $\theta_{U,V}$, $\theta_{V,W}$,
$\theta_{U,W}$ to the inverse images of $U\cap V\cap W$ in $X_V$, $X_W$,
$X_W$ respectively, it follows from the preceding definitions that
$\theta'_{U,V}\circ\theta'_{V,W}=\theta'_{U,W}$. The existence of $X$
satisfying the stated properties therefore follows from
(\hyperref[I.2.3.1]{I, 2.3.1}); its uniqueness up to $Y$-isomorphism is
trivial, taking into account (\hyperref[II.3.1.2.1]{3.1.2.1}).
\end{proof}

\begin{env}[3.1.3]
\label{II.3.1.3}
We will say that the prescheme $X$ defined in
(\hyperref[II.3.1.2]{3.1.2}) is the \emph{homogeneous spectrum} of the
quasi-coherent graded $\mathcal{O}_Y$-algebra $\mathcal{S}$, and we will
denote it by $\operatorname{Proj}(\mathcal{S})$. It is immediate that
$\operatorname{Proj}(\mathcal{S})$ is \emph{separated over $Y$}
((\hyperref[II.2.4.2]{2.4.2}) and
(\hyperref[I.5.5.5]{I, 5.5.5})); if $\mathcal{S}$ is an
$\mathcal{O}_Y$-algebra \emph{of finite type}
(\hyperref[I.9.6.2]{I, 9.6.2}), $\operatorname{Proj}(\mathcal{S})$ is
\emph{of finite type} over $Y$ ((\hyperref[II.2.7.1]{2.7.1, (ii)}) and
(\hyperref[I.6.3.1]{I, 6.3.1})).

If $f$ is the structural morphism $X\to Y$, it is immediate that, for every
prescheme induced by $Y$ on an open subset $U$ of $Y$, $f^{-1}(U)$ is
identified with the homogeneous spectrum $\operatorname{Proj}(\mathcal{S}|U)$.
\end{env}

\begin{proposition}[3.1.4]
\label{II.3.1.4}
Let $f\in\Gamma(Y,\mathcal{S}_d)$ for $d>0$. There exists an open subset
$X_f$ of the underlying space of $X=\operatorname{Proj}(\mathcal{S})$ having
the following property: for every affine open subset $U$ of $Y$,

\oldpage[II]{51}

$X_f\cap\varphi^{-1}(U)=D_+(f|U)$ in $\varphi^{-1}(U)$ identified with
$X_U=\operatorname{Proj}(\Gamma(U,\mathcal{S}))$ ($\varphi$ denoting the
structural morphism $X\to Y$). Moreover, the prescheme induced on $X_f$ is
affine over $Y$ and is canonically isomorphic to
$\operatorname{Spec}(\mathcal{S}^{(d)}/(f-1)\mathcal{S}^{(d)})$
(\hyperref[II.1.3.1]{1.3.1}).
\end{proposition}

\begin{proof}
We have $f|U\in\Gamma(U,\mathcal{S}_d)=(\Gamma(U,\mathcal{S}))_d$. If $U$,
$U'$ are two affine open subsets of $Y$ such that $U'\subset U$, $f|U'$ is
the image of $f|U$ under the restriction homomorphism
\[
  \Gamma(U,\mathcal{S})\longrightarrow\Gamma(U',\mathcal{S}),
\]
so $D_+(f|U')$ is equal (with the notation of
(\hyperref[II.3.1.1]{3.1.1})) to the prescheme induced on the inverse image
$\rho_{U',U}^{-1}(D_+(f|U))$ in $X_{U'}$
(\hyperref[II.2.8.1]{2.8.1}); whence the first assertion. Moreover, the
prescheme induced by $X_U$ on $D_+(f|U)$ is canonically identified with
$\operatorname{Spec}((\Gamma(U,\mathcal{S}))_{(f|U)})$, these identifications
being compatible with the restriction homomorphisms
(\hyperref[II.2.8.1]{2.8.1}); the second assertion then follows from
(\hyperref[II.2.2.5]{2.2.5}) and the commutativity of the diagram
(\hyperref[II.2.8.2.1]{2.8.2.1}).
\end{proof}

We will also say that $X_f$ (as an open subset of the underlying space $X$)
is the set of $x\in X$ where $f$ \emph{does not vanish}.

\begin{corollary}[3.1.5]
\label{II.3.1.5}
If $f\in\Gamma(Y,\mathcal{S}_d)$, $g\in\Gamma(Y,\mathcal{S}_e)$, we have
\[
\label{II.3.1.5.1}
  X_{fg}=X_f\cap X_g.
\tag{3.1.5.1}
\]
\end{corollary}

\begin{proof}
Indeed, it suffices to consider the intersection of the two sides with a set
$\varphi^{-1}(U)$, where $U$ is an affine open subset of $Y$, and to apply
the formula (\hyperref[II.2.3.3.2]{2.3.3.2}).
\end{proof}

\begin{corollary}[3.1.6]
\label{II.3.1.6}
Let $(f_\alpha)$ be a family of sections of $\mathcal{S}$ over $Y$ such that
$f_\alpha\in\Gamma(Y,\mathcal{S}_{d_\alpha})$; if the sheaf of ideals of
$\mathcal{S}$ generated by this family
(\hyperref[0.5.1.1]{0, 5.1.1}) contains all the $\mathcal{S}_n$ from some
index on, the underlying space $X$ is the union of the $X_{f_\alpha}$.
\end{corollary}

\begin{proof}
Indeed, for every affine open subset $U$ of $Y$, $\varphi^{-1}(U)$ is the
union of the $X_{f_\alpha}\cap\varphi^{-1}(U)$
(\hyperref[II.2.3.14]{2.3.14}).
\end{proof}

\begin{corollary}[3.1.7]
\label{II.3.1.7}
Let $\mathcal{A}$ be a quasi-coherent $\mathcal{O}_Y$-algebra; set
\[
  \mathcal{S}=\mathcal{A}[T]
  =\mathcal{A}\otimes_{\mathbf{Z}}\mathbf{Z}[T]
\]
where $T$ is an indeterminate ($\mathbf{Z}$ and $\mathbf{Z}[T]$ being
regarded as simple sheaves on $Y$). Then
$X=\operatorname{Proj}(\mathcal{S})$ is canonically identified with
$\operatorname{Spec}(\mathcal{A})$. In particular,
$\operatorname{Proj}(\mathcal{O}_Y[T])$ is identified with $Y$.
\end{corollary}

\begin{proof}
Applying (\hyperref[II.3.1.6]{3.1.6}) to the unique section
$f\in\Gamma(Y,\mathcal{S})$ equal to $T$ at every point of $Y$, we see that
$X_f=X$. Moreover, here we have $d=1$, and
$\mathcal{S}^{(1)}/(f-1)\mathcal{S}^{(1)}
=\mathcal{S}/(f-1)\mathcal{S}$ is canonically isomorphic to $\mathcal{A}$,
whence the corollary (\hyperref[II.3.1.4]{3.1.4}).
\end{proof}

Let $g\in\Gamma(Y,\mathcal{O}_Y)$; if we take
$\mathcal{S}=\mathcal{O}_Y[T]$, we then have $g\in\Gamma(Y,\mathcal{S}_0)$; let
\[
  h=gT\in\Gamma(Y,\mathcal{S}_1).
\]
If $X=\operatorname{Proj}(\mathcal{S})$, the canonical identification defined
in (\hyperref[II.3.1.7]{3.1.7}) identifies $X_h$ with the open subset
$Y_g$ of $Y$ (in the sense of (\hyperref[0.5.5.2]{0, 5.5.2})): indeed, we
may restrict to the case where $Y=\operatorname{Spec}(A)$ is affine, and
everything then comes down (taking into account
(\hyperref[II.2.2.5]{2.2.5})) to the fact that the ring of fractions $A_g$
is canonically identified with $A[T]/(gT-1)A[T]$
(\hyperref[0.1.2.3]{0, 1.2.3}).

\begin{proposition}[3.1.8]
\label{II.3.1.8}
Let $\mathcal{S}$ be a quasi-coherent $\mathcal{O}_Y$-algebra graded in
positive degrees.
\begin{enumerate}
  \item[\rm(i)] For every $d>0$, there exists a canonical $Y$-isomorphism of
    $\operatorname{Proj}(\mathcal{S})$ onto
    $\operatorname{Proj}(\mathcal{S}^{(d)})$.

\oldpage[II]{52}

  \item[\rm(ii)] Let $\mathcal{S}'$ be the graded $\mathcal{O}_Y$-algebra that
    is the direct sum of $\mathcal{O}_Y$ and the $\mathcal{S}_n$ ($n>0$);
    then $\operatorname{Proj}(\mathcal{S}')$ and
    $\operatorname{Proj}(\mathcal{S})$ are canonically $Y$-isomorphic.
  \item[\rm(iii)] Let $\mathcal{L}$ be an invertible $\mathcal{O}_Y$-module
    (\hyperref[0.5.4.1]{0, 5.4.1}) and let $\mathcal{S}_{(\mathcal{L})}$
    be the graded $\mathcal{O}_Y$-algebra that is the direct sum of the
    $\mathcal{S}_d\otimes\mathcal{L}^{\otimes d}$ ($d\geq0$); then
    $\operatorname{Proj}(\mathcal{S})$ and
    $\operatorname{Proj}(\mathcal{S}_{(\mathcal{L})})$ are canonically
    $Y$-isomorphic.
\end{enumerate}
\end{proposition}

\begin{proof}
In each of the three cases, it suffices to define the isomorphism locally on
$Y$, the verification of the compatibilities with the operations of
restriction from an open subset to a smaller one being trivial. We may
therefore suppose $Y$ affine, and then (i) follows from
(\hyperref[II.2.4.7]{2.4.7, (i)}) and (ii) from
(\hyperref[II.2.4.8]{2.4.8}). As for (iii), when we suppose moreover that
$\mathcal{L}$ is isomorphic to $\mathcal{O}_Y$ (which is permissible, the
question being local on $Y$), the fact that
$\operatorname{Proj}(\mathcal{S})$ and
$\operatorname{Proj}(\mathcal{S}_{(\mathcal{L})})$ are isomorphic is evident;
to define a \emph{canonical} isomorphism, let
$Y=\operatorname{Spec}(A)$, $\mathcal{S}=\widetilde{S}$, where $S$ is a
graded $A$-algebra, and let $c$ be a generator of the free $A$-module $L$
such that $\mathcal{L}=\widetilde{L}$; then, for every $n>0$,
$x_n\mapsto x_n\otimes c^{\otimes n}$ is an $A$-isomorphism of $S_n$ onto
$S_n\otimes L^{\otimes n}$, and these $A$-isomorphisms define an
$A$-isomorphism of graded algebras
\[
  \varphi_c:S\longrightarrow S_{(L)}
  =\bigoplus_{n\geq0}S_n\otimes L^{\otimes n}.
\]
Now let $f\in S_+$ be homogeneous of degree $d$; for every $x\in S_{nd}$, we
have
\[
  \frac{x\otimes c^{\otimes nd}}{(f\otimes c^{\otimes d})^n}
  =\frac{x\otimes(\varepsilon c)^{\otimes nd}}
  {(f\otimes(\varepsilon c)^{\otimes d})^n}
\]
for every invertible element $\varepsilon\in A$, which shows that the
isomorphism $S_{(f)}\to(S_{(L)})_{(f\otimes c^d)}$ deduced from $\varphi_c$
is \emph{independent} of the generator $c$ of $L$ under consideration, and
completes the proof.
\end{proof}

\begin{env}[3.1.9]
\label{II.3.1.9}
Recall ((\hyperref[0.4.1.3]{0, 4.1.3}) and
(\hyperref[I.1.3.14]{I, 1.3.14})) that, for the quasi-coherent graded
$\mathcal{O}_Y$-algebra $\mathcal{S}$ to be \emph{generated by the
$\mathcal{O}_Y$-module $\mathcal{S}_1$}, it is necessary and sufficient that
there exist a covering $(U_\alpha)$ of $Y$ by affine open subsets such that
the graded algebra $\Gamma(U_\alpha,\mathcal{S})$ over
$\Gamma(U_\alpha,\mathcal{S}_0)$ is generated by the set
$\Gamma(U_\alpha,\mathcal{S}_1)$ of its homogeneous elements of degree~$1$.
For every open subset $V$ of $Y$, $\mathcal{S}|V$ is then generated by the
$(\mathcal{O}_Y|V)$-module $\mathcal{S}_1|V$.
\end{env}

\begin{proposition}[3.1.10]
\label{II.3.1.10}
Suppose that there exists a finite affine open covering $(U_i)$ of $Y$ such
that each of the graded algebras $\Gamma(U_i,\mathcal{S})$ is of finite type
over $\Gamma(U_i,\mathcal{O}_Y)$. Then there exists $d>0$ such that
$\mathcal{S}^{(d)}$ is generated by $\mathcal{S}_d$, $\mathcal{S}_d$ being
an $\mathcal{O}_Y$-module of finite type.
\end{proposition}

\begin{proof}
Indeed, it follows from (\hyperref[II.2.1.6]{2.1.6, (v)}) that for each
$i$ there exists an integer $m_i$ such that
$\Gamma(U_i,\mathcal{S}_{nm_i})=(\Gamma(U_i,\mathcal{S}_{m_i}))^n$ for every
$n>0$; it suffices to take for $d$ a common multiple of the $m_i$, taking
into account (\hyperref[II.2.1.6]{2.1.6, (i)}).
\end{proof}

\begin{corollary}[3.1.11]
\label{II.3.1.11}
Under the hypotheses of (\hyperref[II.3.1.10]{3.1.10}),
$\operatorname{Proj}(\mathcal{S})$ is $Y$-isomorphic to a homogeneous spectrum
$\operatorname{Proj}(\mathcal{S}')$, where $\mathcal{S}'$ is a graded
$\mathcal{O}_Y$-algebra generated by $\mathcal{S}'_1$,
$\mathcal{S}'_1$ being assumed to be an $\mathcal{O}_Y$-module of finite type.
\end{corollary}

\begin{proof}
Indeed, it suffices to take $\mathcal{S}'=\mathcal{S}^{(d)}$, $d$ being
determined by the property of (\hyperref[II.3.1.10]{3.1.10}), and to
apply (\hyperref[II.3.1.8]{3.1.8, (i)}).
\end{proof}

\begin{env}[3.1.12]
\label{II.3.1.12}
If $\mathcal{S}$ is a quasi-coherent $\mathcal{O}_Y$-algebra graded in
positive degrees, we know (\hyperref[I.5.1.1]{I, 5.1.1}) that its
\emph{nilradical} $\mathcal{N}$ is a quasi-coherent $\mathcal{O}_Y$-module;
we will say that $\mathcal{N}_+=\mathcal{N}\cap\mathcal{S}_+$ is the
\emph{nilradical of $\mathcal{S}_+$}; it is a quasi-coherent graded
$\mathcal{S}_0$-module, since we are immediately reduced to the case where
$Y$ is affine, and the proposition then follows from
(\hyperref[II.2.1.10]{2.1.10}). For every $y\in Y$,
$(\mathcal{N}_+)_y$ is then the nilradical of
$(\mathcal{S}_+)_y=(\mathcal{S}_y)_+$
(\hyperref[I.5.1.1]{I, 5.1.1}). We say that the graded
$\mathcal{O}_Y$-algebra $\mathcal{S}$ is \emph{essentially reduced} if
$\mathcal{N}_+=0$, which is equivalent

\oldpage[II]{53}

to saying that $\mathcal{S}_y$ is an essentially reduced graded
$\mathcal{O}_y$-algebra for every $y\in Y$. For every graded
$\mathcal{O}_Y$-algebra $\mathcal{S}$, $\mathcal{S}/\mathcal{N}_+$ is
essentially reduced.

We will say that $\mathcal{S}$ is \emph{integral} if $\mathcal{S}_y$ is an
integral ring for every $y\in Y$ and if moreover
$(\mathcal{S}_y)_+=(\mathcal{S}_+)_y\neq0$ for every $y\in Y$.
\end{env}

\begin{proposition}[3.1.13]
\label{II.3.1.13}
Let $\mathcal{S}$ be an $\mathcal{O}_Y$-algebra graded in positive degrees.
If $X=\operatorname{Proj}(\mathcal{S})$, the $Y$-scheme $X_{\mathrm{red}}$ is
canonically isomorphic to $\operatorname{Proj}(\mathcal{S}/\mathcal{N}_+)$;
in particular, if $\mathcal{S}$ is essentially reduced, $X$ is reduced.
\end{proposition}

\begin{proof}
The fact that $X'=\operatorname{Proj}(\mathcal{S}/\mathcal{N}_+)$ is reduced
follows immediately from (\hyperref[II.2.4.4]{2.4.4, (i)}), the property
being local; moreover, for every affine open subset $U\subset Y$,
${\varphi'}^{-1}(U)$ is equal to $(\varphi^{-1}(U))_{\mathrm{red}}$ (denoting
by $\varphi$ and $\varphi'$ the structural morphisms $X\to Y$,
$X'\to Y$); we verify immediately that the canonical $U$-morphisms
${\varphi'}^{-1}(U)\to\varphi^{-1}(U)$ are compatible with the restriction
operations and consequently define a closed immersion $X'\to X$ that is a
homeomorphism of the underlying spaces; whence the conclusion
(\hyperref[I.5.1.2]{I, 5.1.2}).
\end{proof}

\begin{proposition}[3.1.14]
\label{II.3.1.14}
Let $Y$ be an integral prescheme and $\mathcal{S}$ a quasi-coherent graded
$\mathcal{O}_Y$-algebra such that $\mathcal{S}_0=\mathcal{O}_Y$.
\begin{enumerate}
  \item[\rm(i)] If $\mathcal{S}$ is integral
    (\hyperref[II.3.1.12]{3.1.12}), then
    $X=\operatorname{Proj}(\mathcal{S})$ is integral and the structural
    morphism $\varphi:X\to Y$ is dominant.
  \item[\rm(ii)] Suppose moreover that $\mathcal{S}$ is essentially reduced.
    Then, conversely, if $X$ is integral and $\varphi$ is dominant,
    $\mathcal{S}$ is integral.
\end{enumerate}
\end{proposition}

\begin{proof}
\begin{enumerate}
  \item[\rm(i)] If $(U_\alpha)$ is a basis of $Y$ consisting of non-empty
    affine open subsets, it suffices to prove the proposition when $Y$ is
    replaced by one of the $U_\alpha$ and $\mathcal{S}$ by
    $\mathcal{S}|U_\alpha$: indeed, it will follow on the one hand that the
    underlying spaces $\varphi^{-1}(U_\alpha)$ are irreducible (hence
    non-empty) open subsets of $X$ such that
    $\varphi^{-1}(U_\alpha)\cap\varphi^{-1}(U_\beta)\neq\varnothing$ for every
    pair of indices (since $U_\alpha\cap U_\beta$ contains some $U_\gamma$),
    hence that $X$ is irreducible (\hyperref[0.2.1.4]{0, 2.1.4}); on the
    other hand $X$ will be reduced, since this is a local property, so $X$
    will indeed be integral and $\varphi(X)$ dense in $Y$.

    So suppose that $Y=\operatorname{Spec}(A)$, where $A$ is integral
    (\hyperref[I.5.1.4]{I, 5.1.4}), and $\mathcal{S}=\widetilde{S}$, where
    $S$ is a graded $A$-algebra; the hypothesis is that, for every $y\in Y$,
    $\widetilde{S}_y=S_y$ is an integral graded ring such that
    $(S_y)_+\neq0$. It suffices to prove that $S$ is an \emph{integral} ring,
    since we will then have $S_+\neq0$ and we will be able to apply
    (\hyperref[II.2.4.4]{2.4.4, (ii)}). Now let $f$, $g$ be two elements
    $\neq0$ of $S$ and suppose that $fg=0$; for every $y\in Y$ we will then
    have $(f/1)(g/1)=0$ in $S_y$, so $f/1=0$ or $g/1=0$ by hypothesis.
    Suppose for example that $f/1=0$ in $S_y$; this means that there exists
    $a\in A$ such that $a\notin\mathfrak{j}_y$ and $af=0$; we then have, for
    every $z\in Y$, $(a/1)(f/1)=0$ in the \emph{integral} ring $S_z$, and
    since $a/1\neq0$ (because $A$ is integral), $f/1=0$, which implies $f=0$.
  \item[\rm(ii)] The question being local on $Y$, we may again suppose
    that $Y=\operatorname{Spec}(A)$, with $A$ integral, and
    $\mathcal{S}=\widetilde{S}$. By hypothesis, for every $y\in Y$, $(S_y)_+$
    contains no nilpotent element $\neq0$, and the same is true of
    $(S_0)_y=A_y$ by hypothesis; so $S_y$ is a reduced ring for every
    $y\in Y$, and we conclude first that $S$ itself is reduced
    (\hyperref[I.5.1.1]{I, 5.1.1}). On the other hand, the hypothesis that
    $X$ is integral implies that $S$ is essentially integral
    (\hyperref[II.2.4.4]{2.4.4, (ii)}), and everything then comes down to
    showing that the annihilator $\mathfrak{J}$ of $S_+$ in $A=S_0$ is
    reduced

\oldpage[II]{54}

    to $0$
    (\hyperref[II.2.1.11]{2.1.11}). Otherwise, we would have
    $(S_h)_+=0$ for some $h\neq0$ in $\mathfrak{J}$, hence
    (\hyperref[II.3.1.1]{3.1.1}) $\varphi^{-1}(D(h))=\varnothing$, and
    $\varphi(X)$ would not be dense in $Y$, contrary to the hypothesis
    (since $D(h)\neq\varnothing$, $h$ not being nilpotent).
\end{enumerate}
\end{proof}

\subsection{Sheaf on $\operatorname{Proj}(\mathcal{S})$ associated to a
graded $\mathcal{S}$-module.}
\label{subsection:II.3.2}

\begin{env}[3.2.1]
\label{II.3.2.1}
Let $Y$ be a prescheme, $\mathcal{S}$ a quasi-coherent $\mathcal{O}_Y$-algebra
graded in positive degrees, and $\mathcal{M}$ a quasi-coherent graded
$\mathcal{S}$-module (on $(Y,\mathcal{O}_Y)$ or on the ringed space
$(Y,\mathcal{S})$, which amounts to the same
(\hyperref[I.9.6.1]{I, 9.6.1})). With the notation of
(\hyperref[II.3.1.1]{3.1.1}), denote by $\widetilde{\mathcal{M}}_U$ the
quasi-coherent $\mathcal{O}_{X_U}$-module
$\widetilde{\Gamma(U,\mathcal{M})}$; for $U'\subset U$,
$\Gamma(U',\mathcal{M})$ is canonically identified with
$\Gamma(U,\mathcal{M})\otimes_A A'$ (\hyperref[I.1.6.4]{I, 1.6.4}); hence
we have $\widetilde{\mathcal{M}}_{U'}
=\rho_{U',U}^*(\widetilde{\mathcal{M}}_U)$
(\hyperref[II.2.8.11]{2.8.11}).
\end{env}

\begin{proposition}[3.2.2]
\label{II.3.2.2}
There exists on $\operatorname{Proj}(\mathcal{S})=X$ one and only one
quasi-coherent $\mathcal{O}_X$-module $\widetilde{\mathcal{M}}$ such that,
for every affine open subset $U$ of $Y$, we have
\[
  \eta_U^*(\widetilde{\Gamma(U,\mathcal{M})})
  =\widetilde{\mathcal{M}}|f^{-1}(U)
\]
(denoting by $\eta_U$ the isomorphism
$f^{-1}(U)\xrightarrow{\sim}\operatorname{Proj}(\Gamma(U,\mathcal{S}))$, where
$f$ is the structural morphism $X\to Y$).
\end{proposition}

\begin{proof}
Since $\rho_{U',U}$ is identified with the injection morphism
$f^{-1}(U')\to f^{-1}(U)$ (\hyperref[II.3.1.2.1]{3.1.2.1}), the
proposition follows immediately from the relation
$\widetilde{\mathcal{M}}_{U'}=\rho_{U',U}^*(\widetilde{\mathcal{M}}_U)$ and
the principle of gluing sheaves (\hyperref[0.3.3.1]{0, 3.3.1}).
\end{proof}

We say that $\widetilde{\mathcal{M}}$ is the $\mathcal{O}_X$-module associated
to the quasi-coherent graded $\mathcal{S}$-module $\mathcal{M}$.

\begin{proposition}[3.2.3]
\label{II.3.2.3}
Let $\mathcal{M}$ be a quasi-coherent graded $\mathcal{S}$-module, and let
$f\in\Gamma(Y,\mathcal{S}_d)$ ($d>0$). If $\xi_f$ is the canonical isomorphism
of $X_f$ onto the $Y$-prescheme
$Z_f=\operatorname{Spec}(\mathcal{S}^{(d)}/(f-1)\mathcal{S}^{(d)})$
(\hyperref[II.3.1.4]{3.1.4}), then $(\xi_f)_*(\widetilde{\mathcal{M}}|X_f)$
is the $\mathcal{O}_{Z_f}$-module
$\widetilde{\mathcal{M}^{(d)}/(f-1)\mathcal{M}^{(d)}}$
(\hyperref[II.1.4.3]{1.4.3}).
\end{proposition}

\begin{proof}
The question being local on $Y$, we are immediately reduced to
(\hyperref[II.2.2.5]{2.2.5}), taking into account the commutativity of the
diagram (\hyperref[II.2.8.12.1]{2.8.12.1}).
\end{proof}

\begin{proposition}[3.2.4]
\label{II.3.2.4}
The $\mathcal{O}_X$-module $\widetilde{\mathcal{M}}$ is an exact additive
covariant functor in $\mathcal{M}$, from the category of quasi-coherent
graded $\mathcal{S}$-modules to that of quasi-coherent $\mathcal{O}_X$-modules,
which commutes with inductive limits and with direct sums.
\end{proposition}

\begin{proof}
The question, being local on $Y$, reduces to
(\hyperref[I.1.3.11]{I, 1.3.11} and
\hyperref[I.1.3.9]{I, 1.3.9}) and to
(\hyperref[II.2.5.4]{2.5.4}).
\end{proof}

In particular, if $\mathcal{N}$ is a quasi-coherent graded
sub-$\mathcal{S}$-module of $\mathcal{M}$, $\widetilde{\mathcal{N}}$ is
canonically identified with a quasi-coherent sub-$\mathcal{O}_X$-module of
$\widetilde{\mathcal{M}}$; more particularly, for every quasi-coherent graded
sheaf of ideals $\mathcal{J}$ of $\mathcal{S}$,
$\widetilde{\mathcal{J}}$ is a quasi-coherent sheaf of ideals of
$\mathcal{O}_X$.

If $\mathcal{M}$ is a quasi-coherent graded $\mathcal{S}$-module and
$\mathcal{I}$ a quasi-coherent sheaf of ideals of $\mathcal{O}_Y$,
$\mathcal{I}\mathcal{M}$ is a quasi-coherent graded sub-$\mathcal{S}$-module
of $\mathcal{M}$ and we have
\[
\label{II.3.2.4.1}
  \widetilde{\mathcal{I}\mathcal{M}}
  =\mathcal{I}\cdot\widetilde{\mathcal{M}}
\tag{3.2.4.1}
\]
(the right-hand side having the meaning defined in
(\hyperref[0.4.3.5]{0, 4.3.5})). Indeed, it suffices to verify this
formula when $Y=\operatorname{Spec}(A)$ is affine,
$\mathcal{S}=\widetilde{S}$, where $S$ is a graded $A$-algebra,
$\mathcal{M}=\widetilde{M}$,

\oldpage[II]{55}

where $M$ is a graded $S$-module, and
$\mathcal{I}=\widetilde{\mathfrak{I}}$, where $\mathfrak{I}$ is an ideal of
$A$. For every homogeneous $f$ in $S_+$, the restriction to
$D_+(f)=\operatorname{Spec}(S_{(f)})$ of the left-hand side of
(\hyperref[II.3.2.4.1]{3.2.4.1}) is associated to
$(\mathfrak{I}M)_{(f)}=\mathfrak{I}\cdot M_{(f)}$, and the same is true of
the restriction of the right-hand side, in view of
(\hyperref[I.1.3.13]{I, 1.3.13} and
\hyperref[I.1.6.9]{I, 1.6.9}).

\begin{proposition}[3.2.5]
\label{II.3.2.5}
Let $f\in\Gamma(Y,\mathcal{S}_d)$ ($d>0$). On the open subset $X_f$, the
$(\mathcal{O}_X|X_f)$-module
$\widetilde{\mathcal{S}(nd)}|X_f$ is canonically isomorphic to
$\mathcal{O}_X|X_f$ for every $n\in\mathbf{Z}$. In particular, if the
$\mathcal{O}_Y$-algebra $\mathcal{S}$ is generated by $\mathcal{S}_1$
(\hyperref[II.3.1.9]{3.1.9}), the $\mathcal{O}_X$-modules
$\widetilde{\mathcal{S}(n)}$ are invertible for every $n\in\mathbf{Z}$.
\end{proposition}

\begin{proof}
Indeed, for every affine open subset $U$ of $Y$, we have defined in
(\hyperref[II.2.5.7]{2.5.7}) a canonical isomorphism of
$\widetilde{\mathcal{S}(nd)}|(X_f\cap\varphi^{-1}(U))$ onto
$\mathcal{O}_X|(X_f\cap\varphi^{-1}(U))$, taking into account
(\hyperref[II.3.1.4]{3.1.4}) (where $\varphi$ is the structural morphism
$X\to Y$); it is immediately verified that these isomorphisms are compatible
with the restriction from $U$ to an affine open subset $U'\subset U$, whence
the first assertion. To establish the second, it suffices to remark that if
$\mathcal{S}$ is generated by $\mathcal{S}_1$, there is a covering
$(U_\alpha)$ of $Y$ by affine open subsets such that
$\Gamma(U_\alpha,\mathcal{S})$ is generated by
$\Gamma(U_\alpha,\mathcal{S})_1=\Gamma(U_\alpha,\mathcal{S}_1)$; we are then
reduced to applying the result of (\hyperref[II.2.5.9]{2.5.9}), the
property of being invertible being local.
\end{proof}

We also set, for every $n\in\mathbf{Z}$,
\[
\label{II.3.2.5.1}
  \mathcal{O}_X(n)=\widetilde{\mathcal{S}(n)}
\tag{3.2.5.1}
\]
and, for every $\mathcal{O}_X$-module $\mathcal{F}$,
\[
\label{II.3.2.5.2}
  \mathcal{F}(n)=\mathcal{F}\otimes_{\mathcal{O}_X}\mathcal{O}_X(n).
\tag{3.2.5.2}
\]
It follows immediately from these definitions that, for every open subset
$U$ of $Y$, we have
\[
  \widetilde{(\mathcal{S}|U)(n)}
  =\mathcal{O}_X(n)|f^{-1}(U)
\]
$f$ being the structural morphism $X\to Y$.

\begin{proposition}[3.2.6]
\label{II.3.2.6}
Let $\mathcal{M}$ and $\mathcal{N}$ be two quasi-coherent graded
$\mathcal{S}$-modules; there exist a canonical homomorphism, functorial (in
$\mathcal{M}$ and $\mathcal{N}$),
\[
\label{II.3.2.6.1}
  \lambda:\widetilde{\mathcal{M}}
  \otimes_{\mathcal{O}_X}\widetilde{\mathcal{N}}
  \longrightarrow\widetilde{\mathcal{M}\otimes_{\mathcal{S}}\mathcal{N}}
\tag{3.2.6.1}
\]
and a canonical homomorphism, functorial (in $\mathcal{M}$ and
$\mathcal{N}$),
\[
\label{II.3.2.6.2}
  \mu:\widetilde{\mathcal{H}om_{\mathcal{S}}(\mathcal{M},\mathcal{N})}
  \longrightarrow
  \mathcal{H}om_{\mathcal{O}_X}
  (\widetilde{\mathcal{M}},\widetilde{\mathcal{N}}).
\tag{3.2.6.2}
\]
Moreover, if $\mathcal{S}$ is generated by $\mathcal{S}_1$
(\hyperref[II.3.1.9]{3.1.9}), $\lambda$ is an isomorphism; if in addition
$\mathcal{M}$ admits a finite presentation
(\hyperref[II.3.1.1]{3.1.1}), $\mu$ is an isomorphism.
\end{proposition}

\begin{proof}
The homomorphisms $\lambda$ and $\mu$ were defined in
(\hyperref[II.2.5.11.2]{2.5.11.2}) and
(\hyperref[II.2.5.12.2]{2.5.12.2}) when $Y$ is affine; these definitions,
being local, carry over immediately to the general case considered here,
taking into account (\hyperref[II.2.8.14]{2.8.14}).
\end{proof}

\begin{corollary}[3.2.7]
\label{II.3.2.7}
If $\mathcal{S}$ is generated by $\mathcal{S}_1$, we have, for all
$m$, $n$ in $\mathbf{Z}$,
\[
\label{II.3.2.7.1}
  \mathcal{O}_X(m)\otimes_{\mathcal{O}_X}\mathcal{O}_X(n)
  =\mathcal{O}_X(m+n)
\tag{3.2.7.1}
\]
\[
\label{II.3.2.7.2}
  \mathcal{O}_X(n)=(\mathcal{O}_X(1))^{\otimes n}
\tag{3.2.7.2}
\]
up to canonical isomorphisms.
\end{corollary}

\oldpage[II]{56}

\begin{corollary}[3.2.8]
\label{II.3.2.8}
Suppose that $\mathcal{S}$ is generated by $\mathcal{S}_1$. For every graded
$\mathcal{S}$-module $\mathcal{M}$ and every $n\in\mathbf{Z}$, we have
\[
\label{II.3.2.8.1}
  (\mathcal{M}(n))^{\sim}=\widetilde{\mathcal{M}}(n)
\tag{3.2.8.1}
\]
up to canonical isomorphism.
\end{corollary}

\begin{proof}
This follows from the corresponding properties for $Y$ affine
(\hyperref[II.2.5.14]{2.5.14} and
\hyperref[II.2.5.15]{2.5.15}), as well as from
(\hyperref[II.2.8.11]{2.8.11}).
\end{proof}

\par\medskip\noindent\textit{Remarks (3.2.9). ---\ }
\egalabel{II.3.2.9}{3.2.9}
\begin{enumerate}
  \item[\rm(i)] If $\mathcal{S}=\mathcal{A}[T]$, $\mathcal{A}$ being a
    quasi-coherent $\mathcal{O}_Y$-algebra
    (\hyperref[II.3.1.7]{3.1.7}), we verify immediately that all the
    invertible $\mathcal{O}_X$-modules $\mathcal{O}_X(n)$ are canonically
    isomorphic to $\mathcal{O}_X$.

    Moreover, let $\mathcal{N}$ be a quasi-coherent $\mathcal{A}$-module, and
    set $\mathcal{M}=\mathcal{N}\otimes_{\mathcal{A}}\mathcal{A}[T]$. It then
    follows from (\hyperref[II.3.2.3]{3.2.3}) and
    (\hyperref[II.3.1.7]{3.1.7}) that, under the canonical identification
    of $X=\operatorname{Proj}(\mathcal{A}[T])$ and
    $X'=\operatorname{Spec}(\mathcal{A})$, the $\mathcal{O}_X$-module
    $\widetilde{\mathcal{M}}$ is identified with the $\mathcal{O}_{X'}$-module
    $\widetilde{\mathcal{N}}$ associated to $\mathcal{N}$ (in the sense of
    (\hyperref[II.1.4.3]{1.4.3})).
  \item[\rm(ii)] Let $\mathcal{S}$ be an arbitrary graded
    $\mathcal{O}_Y$-algebra, and $\mathcal{S}'$ the graded
    $\mathcal{O}_Y$-algebra such that
    $\mathcal{S}'_0=\mathcal{O}_Y$, $\mathcal{S}'_n=\mathcal{S}_n$ for every
    $n>0$; the canonical isomorphism of
    $X=\operatorname{Proj}(\mathcal{S})$ onto
    $X'=\operatorname{Proj}(\mathcal{S}')$
    (\hyperref[II.3.1.8]{3.1.8, (ii)}) identifies $\mathcal{O}_X(n)$ and
    $\mathcal{O}_{X'}(n)$ for every $n\in\mathbf{Z}$: this follows from the
    same proposition for the affine case (\hyperref[II.2.5.16]{2.5.16})
    and from the fact that these identifications, for the affine open subsets
    of $Y$, commute with the restriction operations. Likewise, let
    $X^{(d)}=\operatorname{Proj}(\mathcal{S}^{(d)})$; the canonical
    isomorphism of $X$ onto $X^{(d)}$
    (\hyperref[II.3.1.8]{3.1.8, (i)}) identifies $\mathcal{O}_X(nd)$ and
    $\mathcal{O}_{X^{(d)}}(n)$ for every $n\in\mathbf{Z}$.
\end{enumerate}

\begin{proposition}[3.2.10]
\label{II.3.2.10}
Let $\mathcal{L}$ be an invertible $\mathcal{O}_Y$-module and $g$ the
canonical isomorphism
$X_{(\mathcal{L})}=\operatorname{Proj}(\mathcal{S}_{(\mathcal{L})})
\to X=\operatorname{Proj}(\mathcal{S})$
(\hyperref[II.3.1.8]{3.1.8, (iii)}). For every integer $n\in\mathbf{Z}$,
$g_*(\mathcal{O}_{X_{(\mathcal{L})}}(n))$ is canonically isomorphic to
$\mathcal{O}_X(n)\otimes_Y\mathcal{L}^{\otimes n}$.
\end{proposition}

\begin{proof}
Suppose first that $Y$ is affine, with ring $A$, and that
$\mathcal{L}=\widetilde{L}$, where $L$ is a monogenic free $A$-module. With
the notation of the proof of
(\hyperref[II.3.1.8]{3.1.8, (iii)}), we define, for $f\in S_d$, an
isomorphism of $(S(n))_{(f)}\otimes_A L^{\otimes n}$ onto
$(S_{(L)}(n))_{(f\otimes c^{\otimes d})}$ by sending
$(x/f^k)\otimes c^{\otimes n}$, where $x\in S_{kd+n}$, to the element
$(x\otimes c^{\otimes(n+kd)})/(f\otimes c^{\otimes d})^k$; it is immediate
that this isomorphism is independent of the generator $c$ chosen for $L$;
moreover, the isomorphisms thus defined for each $f\in S_+$ are compatible
with the restriction operations $D_+(f)\to D_+(fg)$. Finally, in the
general case, we see without difficulty, by going back to the definitions
(\hyperref[II.3.1.1]{3.1.1}), that the isomorphisms thus defined for
each affine open subset $U$ of $Y$ are compatible with the passage from $U$
to an affine open subset $U'\subset U$.
\end{proof}

\subsection{Graded $\mathcal{S}$-module associated to a sheaf on
$\operatorname{Proj}(\mathcal{S})$.}
\label{subsection:II.3.3}

\emph{Throughout this subsection we suppose that the graded
$\mathcal{O}_Y$-algebra $\mathcal{S}$ is generated by $\mathcal{S}_1$
(\hyperref[II.3.1.9]{3.1.9}).} Recall that, because of
(\hyperref[II.3.1.8]{3.1.8, (i)}), this restriction is not essential,
subject to the finiteness conditions (\hyperref[II.3.1.10]{3.1.10}).

\begin{env}[3.3.1]
\label{II.3.3.1}
Let $p$ be the structural morphism $X=\operatorname{Proj}(\mathcal{S})\to Y$.
For every $\mathcal{O}_X$-module $\mathcal{F}$, we set
\[
\label{II.3.3.1.1}
  \Gamma_*(\mathcal{F})
  =\bigoplus_{n\in\mathbf{Z}}p_*(\mathcal{F}(n))
\tag{3.3.1.1}
\]

\oldpage[II]{57}

and in particular
\[
\label{II.3.3.1.2}
  \Gamma_*(\mathcal{O}_X)
  =\bigoplus_{n\in\mathbf{Z}}p_*(\mathcal{O}_X(n)).
\tag{3.3.1.2}
\]
We know (\hyperref[0.4.2.2]{0, 4.2.2}) that there exists a canonical
homomorphism
\[
  p_*(\mathcal{F})\otimes_{\mathcal{O}_Y}p_*(\mathcal{G})
  \longrightarrow
  p_*(\mathcal{F}\otimes_{\mathcal{O}_X}\mathcal{G})
\]
for two $\mathcal{O}_X$-modules $\mathcal{F}$ and $\mathcal{G}$; we therefore
deduce from (\hyperref[II.3.2.7.1]{3.2.7.1}) that
$\Gamma_*(\mathcal{O}_X)$ is endowed with a structure of
\emph{graded $\mathcal{O}_Y$-algebra}, and
(\hyperref[II.3.2.5.2]{3.2.5.2}) likewise defines on
$\Gamma_*(\mathcal{F})$ a structure of \emph{graded module} over
$\Gamma_*(\mathcal{O}_X)$.

By virtue of (\hyperref[II.3.2.5]{3.2.5}) and of the left exactness of the
functor $p_*$ (\hyperref[0.4.2.1]{0, 4.2.1}),
$\Gamma_*(\mathcal{F})$ is a covariant additive and \emph{left exact}
functor in $\mathcal{F}$ on the category of $\mathcal{O}_X$-modules, with
values in the category of graded $\mathcal{O}_Y$-modules (where the morphisms
are the homomorphisms of degree $0$). In particular, if $\mathcal{J}$ is a
sheaf of ideals in $\mathcal{O}_X$, $\Gamma_*(\mathcal{J})$ is identified
with a \emph{graded sheaf of ideals} of $\Gamma_*(\mathcal{O}_X)$.
\end{env}

\begin{env}[3.3.2]
\label{II.3.3.2}
Let $\mathcal{M}$ be a quasi-coherent graded $\mathcal{S}$-module. For every
affine open subset $U$ of $Y$, we have defined in
(\hyperref[II.2.6.2]{2.6.2}) a homomorphism of abelian groups
\[
  \alpha_{0,U}:\Gamma(U,\mathcal{M}_0)
  \longrightarrow\Gamma(p^{-1}(U),\widetilde{\mathcal{M}}).
\]
It is immediate that these homomorphisms commute with the restriction
operations (\hyperref[II.2.8.13.1]{2.8.13.1}) and therefore define
(without using the hypothesis that $\mathcal{S}$ is generated by
$\mathcal{S}_1$) a homomorphism of sheaves of abelian groups
\[
\label{II.3.3.2.1}
  \alpha_0:\mathcal{M}_0\longrightarrow p_*(\widetilde{\mathcal{M}}).
\tag{3.3.2.1}
\]
Applying this result to each of the
$\mathcal{M}_n=(\mathcal{M}(n))_0$, and taking into account
(\hyperref[II.3.2.8.1]{3.2.8.1}), we define a homomorphism of sheaves of
abelian groups
\[
\label{II.3.3.2.2}
  \alpha_n:\mathcal{M}_n\longrightarrow
  p_*(\widetilde{\mathcal{M}}(n))
  \qquad\text{for every }n\in\mathbf{Z},
\tag{3.3.2.2}
\]
whence a functorial homomorphism (of degree $0$) of graded sheaves of abelian
groups
\[
\label{II.3.3.2.3}
  \alpha:\mathcal{M}\longrightarrow\Gamma_*(\widetilde{\mathcal{M}})
\tag{3.3.2.3}
\]
(also denoted $\alpha_{\mathcal{M}}$).

Taking in particular $\mathcal{M}=\mathcal{S}$, we verify that
$\alpha:\mathcal{S}\to\Gamma_*(\mathcal{O}_X)$ is a homomorphism of graded
$\mathcal{O}_Y$-algebras and that
(\hyperref[II.3.3.2.3]{3.3.2.3}) is a di-homomorphism of graded modules,
relative to this homomorphism of graded algebras.

Note also that to each of the $\alpha_n$ there corresponds
(\hyperref[0.4.4.3]{0, 4.4.3}) a canonical homomorphism of
$\mathcal{O}_X$-modules
\[
\label{II.3.3.2.4}
  \alpha_n^\sharp:p^*(\mathcal{M}_n)
  \longrightarrow\widetilde{\mathcal{M}}(n).
\tag{3.3.2.4}
\]
We verify without difficulty that this homomorphism is none other than the
one that corresponds functorially (\hyperref[II.3.2.4]{3.2.4}) to the
canonical homomorphism (of degree $0$) of graded $\mathcal{O}_Y$-modules
\[
\label{II.3.3.2.5}
  \mathcal{M}_n\otimes_{\mathcal{O}_Y}\mathcal{S}
  \longrightarrow\mathcal{M}(n)
\tag{3.3.2.5}
\]

\oldpage[II]{58}

where the grading of the right-hand side comes naturally from that of
$\mathcal{M}$. Indeed, we may restrict to the case where
$Y=\operatorname{Spec}(A)$ is affine, $\mathcal{M}=\widetilde{M}$ and
$\mathcal{S}=\widetilde{S}$, the graded $A$-algebra $S$ being generated by
$S_1$, so that, as $f$ runs through $S_1$, the $D_+(f)$ form a covering of
$X$. Going back to the definitions
(\hyperref[II.2.6.2]{2.6.2}), we then see, taking into account
(\hyperref[I.1.6.7]{I, 1.6.7}), that the restriction to $D_+(f)$ of
the homomorphism (\hyperref[II.3.3.2.4]{3.3.2.4}) corresponds
(\hyperref[I.1.3.8]{I, 1.3.8}) to the homomorphism of $S_{(f)}$-modules
$M_n\otimes_A S_{(f)}\to(S(n))_{(f)}$ which sends $x\otimes1$ ($x\in M_n$)
to $x/1$; this proves our assertion.
\end{env}

\begin{proposition}[3.3.3]
\label{II.3.3.3}
For every section $f\in\Gamma(Y,\mathcal{S}_d)$ ($d>0$), $X_f$ is identical
to the set of points of $X$ where $\alpha_d(f)$ (regarded as a section of
$\mathcal{O}_X(d)$) does not vanish
(\hyperref[0.5.5.2]{0, 5.5.2}).
\end{proposition}

\begin{proof}
($\alpha_d(f)$ is a section of $p_*(\mathcal{O}_X(d))$ over $Y$, but by
definition such a section is also a section of $\mathcal{O}_X(d)$ over $X$
(\hyperref[0.4.2.1]{0, 4.2.1})). The definition of $X_f$
(\hyperref[II.3.1.4]{3.1.4}) reduces us to the case where $Y$ is affine,
which was treated in (\hyperref[II.2.6.3]{2.6.3}).
\end{proof}

\begin{env}[3.3.4]
\label{II.3.3.4}
From now on we will suppose, in addition to the hypothesis made at the
beginning of this subsection, that for every quasi-coherent
$\mathcal{O}_X$-module $\mathcal{F}$ the $p_*(\mathcal{F}(n))$ are
\emph{quasi-coherent} on $Y$, and consequently
$\Gamma_*(\mathcal{F})=\bigoplus_{n\in\mathbf{Z}}p_*(\mathcal{F}(n))$ is
also a quasi-coherent $\mathcal{O}_Y$-module
(\hyperref[I.1.4.1]{I, 1.4.1} and
\hyperref[I.1.3.9]{I, 1.3.9}); this will always be the case in particular
if $X$ is \emph{of finite type} over $Y$
(\hyperref[I.9.2.2]{I, 9.2.2}). We conclude that
$(\Gamma_*(\mathcal{F}))^{\sim}$ is defined and is a quasi-coherent
$\mathcal{O}_X$-module. For every affine open subset $U$ of $Y$, we
have (\hyperref[I.1.3.9]{I, 1.3.9} and
\hyperref[II.2.5.4]{2.5.4})
\begin{align*}
  \left(\Gamma\left(U,\bigoplus_{n\in\mathbf{Z}}
  p_*(\mathcal{F}(n))\right)\right)^{\sim}
  &=\bigoplus_{n\in\mathbf{Z}}
  \bigl(\Gamma(U,p_*(\mathcal{F}(n)))\bigr)^{\sim}\\
  &=\bigoplus_{n\in\mathbf{Z}}
  \bigl(\Gamma(p^{-1}(U),\mathcal{F}(n))\bigr)^{\sim}\\
  &=\left(\bigoplus_{n\in\mathbf{Z}}
  \Gamma(p^{-1}(U),\mathcal{F}(n))\right)^{\sim}\\
  &=\bigl(\Gamma_*(\mathcal{F}|p^{-1}(U))\bigr)^{\sim}
\end{align*}
and consequently (\hyperref[II.2.6.4]{2.6.4}) we have a canonical
homomorphism
\[
  \beta_U:\left(\Gamma\left(U,\bigoplus_{n\in\mathbf{Z}}
  p_*(\mathcal{F}(n))\right)\right)^{\sim}
  \longrightarrow\mathcal{F}|p^{-1}(U).
\]

Moreover, the commutativity of
(\hyperref[II.2.8.13.2]{2.8.13.2}) shows that these homomorphisms commute
with the restriction operations on $Y$; we therefore deduce a canonical
functorial homomorphism
\[
\label{II.3.3.4.1}
  \beta:(\Gamma_*(\mathcal{F}))^{\sim}\longrightarrow\mathcal{F}
\tag{3.3.4.1}
\]
(also denoted $\beta_{\mathcal{F}}$) for quasi-coherent
$\mathcal{O}_X$-modules.
\end{env}

\begin{proposition}[3.3.5]
\label{II.3.3.5}
Let $\mathcal{M}$ be a quasi-coherent graded $\mathcal{S}$-module and
$\mathcal{F}$ a quasi-coherent $\mathcal{O}_X$-module; the composite
homomorphisms
\[
\label{II.3.3.5.1}
  \widetilde{\mathcal{M}}
  \xrightarrow{\widetilde{\alpha}}
  (\Gamma_*(\widetilde{\mathcal{M}}))^{\sim}
  \xrightarrow{\beta}\widetilde{\mathcal{M}}
\tag{3.3.5.1}
\]
\[
\label{II.3.3.5.2}
  \Gamma_*(\mathcal{F})\xrightarrow{\alpha}
  \Gamma_*((\Gamma_*(\mathcal{F}))^{\sim})
  \xrightarrow{\Gamma_*(\beta)}\Gamma_*(\mathcal{F})
\tag{3.3.5.2}
\]
are the identity isomorphisms.
\end{proposition}

\begin{proof}
Indeed, the question is local on $Y$, and we are reduced to
(\hyperref[II.2.6.5]{2.6.5}).
\end{proof}


\subsection{Finiteness conditions.}
\label{subsection:II.3.4}

\oldpage[II]{59}

\begin{proposition}[3.4.1]
\label{II.3.4.1}
Let $Y$ be a prescheme, and $\mathcal{S}$ a quasi-coherent $\mathcal{O}_Y$-algebra
generated by $\mathcal{S}_1$
(\hyperref[II.3.1.9]{3.1.9}); suppose moreover that $\mathcal{S}_1$ is of
finite type. Then $X=\operatorname{Proj}(\mathcal{S})$ is of finite type over
$Y$.
\end{proposition}

\begin{proof}
Indeed, the question being local on $Y$, we may suppose $Y$ affine
with ring $A$; we then have $\mathcal{S}=\widetilde{S}$, where
$S=\Gamma(Y,\mathcal{S})$, and by hypothesis $S$ is an $A$-algebra generated
by $S_1=\Gamma(Y,\mathcal{S}_1)$, where we may moreover suppose that $S_1$ is
an $A$-module of finite type
(\hyperref[I.1.3.9]{I, 1.3.9} and
\hyperref[I.1.3.12]{I, 1.3.12}). Consequently $S$ is a graded $A$-algebra
of finite type, and we are reduced to
(\hyperref[II.2.7.1]{2.7.1, (ii)}).
\end{proof}

\begin{env}[3.4.2]
\label{II.3.4.2}
Let $\mathcal{S}$ be a quasi-coherent graded $\mathcal{O}_Y$-algebra; for
a quasi-coherent graded $\mathcal{S}$-module $\mathcal{M}$, we will consider
the following finiteness conditions:
\begin{enumerate}
  \item[(\textbf{TF})] \emph{There exists an integer $n$ such that the
    $\mathcal{S}$-module $\bigoplus_{k\geq n}\mathcal{M}_k$ is of finite type.}
  \item[(\textbf{TN})] \emph{There exists an integer $n$ such that
    $\mathcal{M}_k=0$ for $k\geq n$.}
\end{enumerate}
If $\mathcal{M}$ satisfies (\textbf{TN}), we have
$\widetilde{\mathcal{M}}=0$, the question being local on $Y$
(\hyperref[II.2.7.2]{2.7.2}).

Let $\mathcal{M}$, $\mathcal{N}$ be two quasi-coherent graded
$\mathcal{S}$-modules; we say that a homomorphism
$u:\mathcal{M}\to\mathcal{N}$ of degree $0$ is
(\textbf{TN})-\emph{injective} (resp. (\textbf{TN})-\emph{surjective},
(\textbf{TN})-\emph{bijective}) if there exists an integer $n$ such that
$u_k:\mathcal{M}_k\to\mathcal{N}_k$ is injective (resp. surjective, bijective)
for $k\geq n$; then
$\widetilde{u}:\widetilde{\mathcal{M}}\to\widetilde{\mathcal{N}}$ is injective
(resp. surjective, bijective) by virtue of
(\hyperref[II.2.7.2]{2.7.2}), the question being local on $Y$, and taking
into account (\hyperref[I.1.3.9]{I, 1.3.9}); when $u$ is
(\textbf{TN})-bijective, we also say that $u$ is a
(\textbf{TN})-\emph{isomorphism}.
\end{env}

\begin{proposition}[3.4.3]
\label{II.3.4.3}
Let $Y$ be a prescheme, $\mathcal{S}$ a quasi-coherent graded
$\mathcal{O}_Y$-algebra generated by $\mathcal{S}_1$, $\mathcal{S}_1$ being
supposed of finite type. Let $\mathcal{M}$ be a quasi-coherent graded
$\mathcal{S}$-module.
\begin{enumerate}
  \item[\rm(i)] If $\mathcal{M}$ satisfies condition (\textbf{TF}),
    $\widetilde{\mathcal{M}}$ is of finite type.
  \item[\rm(ii)] Suppose that $\mathcal{M}$ satisfies (\textbf{TF}); for
    $\widetilde{\mathcal{M}}=0$, it is necessary and sufficient that
    $\mathcal{M}$ satisfy (\textbf{TN}).
\end{enumerate}
\end{proposition}

\begin{proof}
The questions being local on $Y$, we are reduced to the case where $Y$ is
affine with ring $A$, $\mathcal{S}=\widetilde{S}$, where $S$ is a graded
$A$-algebra such that the ideal $S_+$ is of finite type, and
$\mathcal{M}=\widetilde{M}$ where $M$ is a graded $S$-module; the proposition
then follows from (\hyperref[II.2.7.3]{2.7.3}).
\end{proof}

\begin{theorem}[3.4.4]
\label{II.3.4.4}
Let $Y$ be a prescheme, $\mathcal{S}$ a quasi-coherent graded
$\mathcal{O}_Y$-algebra generated by $\mathcal{S}_1$, $\mathcal{S}_1$ being
supposed of finite type; let $X=\operatorname{Proj}(\mathcal{S})$. For every
quasi-coherent $\mathcal{O}_X$-module $\mathcal{F}$, the canonical homomorphism
$\beta$ (\hyperref[II.3.3.4]{3.3.4}) is an isomorphism.
\end{theorem}

\begin{proof}
Note first that $\beta$ is defined by virtue of
(\hyperref[II.3.4.1]{3.4.1}). To see that $\beta$ is an isomorphism,
we are reduced to the case where $Y$ is affine with ring $A$,
$\mathcal{S}=\widetilde{S}$, where $S$ is a graded $A$-algebra generated by
$S_1$ and $S_1$ is an $A$-module of finite type. It then suffices to apply
(\hyperref[II.2.7.5]{2.7.5}).
\end{proof}

\begin{corollary}[3.4.5]
\label{II.3.4.5}
Under the hypotheses of (\hyperref[II.3.4.4]{3.4.4}), every
quasi-coherent $\mathcal{O}_X$-module $\mathcal{F}$ is isomorphic to an
$\mathcal{O}_X$-module of the form $\widetilde{\mathcal{M}}$, where
$\mathcal{M}$ is a quasi-coherent graded $\mathcal{S}$-module. If moreover
$\mathcal{F}$ is of finite type, and if we suppose that $Y$ is a quasi-compact
scheme, or that the underlying space of $Y$ is Noetherian, we may
suppose $\mathcal{M}$ of finite type.
\end{corollary}

\oldpage[II]{60}

\begin{proof}
The first assertion follows immediately from
(\hyperref[II.3.4.4]{3.4.4}) by taking
$\mathcal{M}=\Gamma_*(\mathcal{F})$. To establish the second, it will suffice
to establish that $\mathcal{M}$ is the inductive limit of its graded
sub-$\mathcal{S}$-modules $\mathcal{N}_\lambda$ of finite type: indeed, it
will follow that $\widetilde{\mathcal{M}}$ is the inductive limit of the
$\widetilde{\mathcal{N}}_\lambda$
(\hyperref[II.3.2.4]{3.2.4}), hence $\mathcal{F}$ is the inductive limit of
the $\beta(\widetilde{\mathcal{N}}_\lambda)$; since $X$ is quasi-compact
((\hyperref[II.3.4.1]{3.4.1}) and
\hyperref[I.6.3.1]{I, 6.3.1}) and $\mathcal{F}$ is of finite type,
$\mathcal{F}$ will necessarily be equal to one of the
$\beta(\widetilde{\mathcal{N}}_\lambda)$
(\hyperref[0.5.2.3]{0, 5.2.3}).

To define the $\mathcal{N}_\lambda$ having $\mathcal{M}$ as inductive limit,
it suffices to consider, for each $n\in\mathbf{Z}$, the quasi-coherent
$\mathcal{O}_Y$-module $\mathcal{M}_n$, which is the inductive limit of its
sub-$\mathcal{O}_Y$-modules $\mathcal{M}_n^{(\mu_n)}$ of finite type, because
of the hypotheses on $Y$
(\hyperref[I.9.4.9]{I, 9.4.9}); it is immediate that
$\mathcal{P}_{\mu_n}=\mathcal{S}\cdot\mathcal{M}_n^{(\mu_n)}$ is a
graded $\mathcal{S}$-module of finite type, and we verify at once that, by
taking for $\mathcal{N}_\lambda$ the finite sums of $\mathcal{S}$-modules of
the form $\mathcal{P}_{\mu_n}$, we do answer the question.
\end{proof}

\begin{corollary}[3.4.6]
\label{II.3.4.6}
Suppose that the hypotheses of (\hyperref[II.3.4.4]{3.4.4}) are
satisfied, and moreover that the underlying space $Y$ is quasi-compact; let
$\mathcal{F}$ be a quasi-coherent $\mathcal{O}_X$-module of finite type. There
exists $n_0$ such that, for $n\geq n_0$, the canonical homomorphism
$\sigma:p^*(p_*(\mathcal{F}(n)))\to\mathcal{F}(n)$
(\hyperref[0.4.4.3]{0, 4.4.3}) is surjective.
\end{corollary}

\begin{proof}
Indeed, for every $y\in Y$, let $U$ be an affine open neighbourhood of $y$ in
$Y$. There exists an integer $n_0(U)$ such that, for $n\geq n_0(U)$,
$\mathcal{F}(n)|p^{-1}(U)$ is generated by finitely many of its sections
over $p^{-1}(U)$ (\hyperref[II.2.7.9]{2.7.9}); but these sections
are canonical images of sections of $p^*(p_*(\mathcal{F}(n)))$ over
$p^{-1}(U)$ (\hyperref[0.3.7.1]{0, 3.7.1} and
\hyperref[0.4.4.3]{0, 4.4.3}), hence
$\mathcal{F}(n)|p^{-1}(U)$ is equal to the canonical image of
$p^*(p_*(\mathcal{F}(n)))|p^{-1}(U)$. Finally, since $Y$ is quasi-compact,
there exists a finite covering of $Y$ by affine open subsets $U_i$, and by
taking for $n_0$ the largest of the $n_0(U_i)$, we complete the proof.
\end{proof}

\par\medskip\noindent\textit{Remarks (3.4.7). ---\ }
\egalabel{II.3.4.7}{3.4.7}
If $p=(\psi,\theta):X\to Y$ is a morphism of ringed spaces and
$\mathcal{F}$ an $\mathcal{O}_X$-module, the fact that the canonical homomorphism
$\sigma:p^*(p_*(\mathcal{F}))\to\mathcal{F}$ is surjective can be stated
explicitly as follows (\hyperref[0.4.4.1]{0, 4.4.1}): for every $x\in X$
and every section $s$ of $\mathcal{F}$ over an open neighbourhood $V$ of $x$,
there exist an open neighbourhood $U$ of $p(x)$ in $Y$, finitely many
sections $t_i$ ($1\leq i\leq m$) of $\mathcal{F}$ over $p^{-1}(U)$, a
neighbourhood $W\subset V\cap p^{-1}(U)$ of $x$ and sections $a_i$
($1\leq i\leq m$) of $\mathcal{O}_X$ over $W$ such that
\[
  s|W=\sum_i a_i\cdot(t_i|W).
\]
When $Y$ is an \emph{affine scheme} and $p_*(\mathcal{F})$ is
\emph{quasi-coherent}, this condition is equivalent to the fact that
$\mathcal{F}$ is \emph{generated by its sections over $X$}
(\hyperref[0.5.5.1]{0, 5.5.1}): indeed, if
$Y=\operatorname{Spec}(A)$, we may suppose that $U=D(f)$ with $f\in A$; there
exist an integer $n>0$ and sections $s_i$ of $\mathcal{F}$ over
$X$ such that $t_i$ is the restriction to $p^{-1}(U)$ of $s_ig^n$, with
$g=\theta(f)$ (by applying (\hyperref[I.1.4.1]{I, 1.4.1}) to
$p_*(\mathcal{F})$); since $g$ is invertible over $p^{-1}(U)$, we
therefore have
\[
  s|W=\sum_i b_i\cdot(s_i|W)
\]
with $b_i=a_i(g|W)^{-n}$, whence our assertion. When $Y$ is affine,
Corollary (\hyperref[II.3.4.6]{3.4.6}) therefore recovers
(\hyperref[II.2.7.9]{2.7.9}).

\oldpage[II]{61}

We conclude that when $Y$ is an arbitrary prescheme, the following three
conditions, for a quasi-coherent $\mathcal{O}_X$-module
$\mathcal{F}$ such that $p_*(\mathcal{F})$ is \emph{quasi-coherent}, are
equivalent:
\begin{enumerate}
  \item[a)] \emph{The canonical homomorphism
    $\sigma:p^*(p_*(\mathcal{F}))\to\mathcal{F}$ is surjective.}
  \item[b)] \emph{There exist a quasi-coherent $\mathcal{O}_Y$-module
    $\mathcal{G}$ and a surjective homomorphism
    $p^*(\mathcal{G})\to\mathcal{F}$.}
  \item[c)] \emph{For every affine open subset $U$ of $Y$,
    $\mathcal{F}|p^{-1}(U)$ is generated by its sections over
    $p^{-1}(U)$.}
\end{enumerate}
We have indeed just proved the equivalence of \emph{a)} and \emph{c)}. On the
other hand, it is clear that \emph{a)} implies \emph{b)}, $p_*(\mathcal{F})$
being quasi-coherent by hypothesis. Conversely, every homomorphism
$u:p^*(\mathcal{G})\to\mathcal{F}$ factors as
$p^*(\mathcal{G})\to p^*(p_*(\mathcal{F}))\xrightarrow{\sigma}\mathcal{F}$
(\hyperref[0.3.5.4.4]{0, 3.5.4.4}), hence if $u$ is surjective so is
$\sigma$, which proves that \emph{b)} implies \emph{a)}.

\begin{corollary}[3.4.8]
\label{II.3.4.8}
Suppose that the hypotheses of (\hyperref[II.3.4.4]{3.4.4}) are satisfied,
and suppose moreover that $Y$ is a quasi-compact scheme or that the
underlying space of $Y$ is Noetherian. Let $\mathcal{F}$ be a
quasi-coherent $\mathcal{O}_X$-module of finite type; there then exists an
integer $n_0$ such that, for $n\geq n_0$, $\mathcal{F}$ is isomorphic to a
quotient of an $\mathcal{O}_X$-module of the form $(p^*(\mathcal{G}))(-n)$,
where $\mathcal{G}$ is a quasi-coherent $\mathcal{O}_Y$-module of finite type
(depending on $n$).
\end{corollary}

\begin{proof}
Since the structural morphism $X\to Y$ is separated and of finite type,
$p_*(\mathcal{F}(n))$ is quasi-coherent
(\hyperref[I.9.2.2]{I, 9.2.2, b)}), hence is the inductive limit of its
quasi-coherent sub-$\mathcal{O}_Y$-modules of finite type, by virtue of
the hypothesis on $Y$ (\hyperref[I.9.4.9]{I, 9.4.9}). We deduce from this,
by (\hyperref[II.3.4.6]{3.4.6}),
(\hyperref[0.4.3.2]{0, 4.3.2}) and
(\hyperref[0.5.2.3]{0, 5.2.3}), that $\mathcal{F}(n)$ is the canonical
image of an $\mathcal{O}_X$-module of the form $p^*(\mathcal{G})$, where
$\mathcal{G}$ is a quasi-coherent sub-$\mathcal{O}_Y$-module of finite type of
$p_*(\mathcal{F}(n))$; the corollary then follows from
(\hyperref[II.3.2.5.2]{3.2.5.2}) and
(\hyperref[II.3.2.7.1]{3.2.7.1}).
\end{proof}

\subsection{Functorial behaviour.}
\label{subsection:II.3.5}

\begin{env}[3.5.1]
\label{II.3.5.1}
Let $Y$ be a prescheme, and $\mathcal{S}$, $\mathcal{S}'$ two quasi-coherent
$\mathcal{O}_Y$-algebras graded in positive degrees; set
$X=\operatorname{Proj}(\mathcal{S})$, $X'=\operatorname{Proj}(\mathcal{S}')$,
and let $p$, $p'$ be the structural morphisms of $X$ and $X'$ to $Y$. Let
$\varphi:\mathcal{S}'\to\mathcal{S}$ be an $\mathcal{O}_Y$-homomorphism
of graded algebras. For every affine open subset $U$ of $Y$, set
$S_U=\Gamma(U,\mathcal{S})$, $S'_U=\Gamma(U,\mathcal{S}')$;
the homomorphism $\varphi$ defines a homomorphism
$\varphi_U:S'_U\to S_U$ of graded $A_U$-algebras, where we set
$A_U=\Gamma(U,\mathcal{O}_Y)$. To it there correspond in $p^{-1}(U)$ an open
set $G(\varphi_U)$ and a morphism
$\Phi_U:G(\varphi_U)\to p'^{-1}(U)$
(\hyperref[II.2.8.1]{2.8.1}). Moreover, if $V\subset U$ is an affine open
subset, the diagram
\[
\label{II.3.5.1.1}
  \xymatrix{
    S'_U\ar[r]^{\varphi_U}\ar[d]
    & S_U\ar[d]\\
    S'_V\ar[r]_{\varphi_V}
    & S_V
  }
\tag{3.5.1.1}
\]
is commutative, and we verify at once, from the definitions
(\hyperref[II.2.8.1]{2.8.1}), that we have
\[
  G(\varphi_V)=G(\varphi_U)\cap p^{-1}(V)
\]
and that $\Phi_V$ is the restriction of $\Phi_U$ to $G(\varphi_V)$. We have
thus defined an open subset $G(\varphi)$ of $X$ such that
$G(\varphi)\cap p^{-1}(U)=G(\varphi_U)$ for every affine open subset $U\subset Y$,
and an affine $Y$-morphism

\oldpage[II]{62}

$\Phi:G(\varphi)\to X'$, which is said to be \emph{associated} to $\varphi$ and
which we will denote by $\operatorname{Proj}(\varphi)$. When, for every
$y\in Y$, there is an affine neighbourhood $U$ of $y$ such that the
$\Gamma(U,\mathcal{O}_Y)$-module $\Gamma(U,\mathcal{S}_+)$ is generated by
$\varphi(\Gamma(U,\mathcal{S}'_+))$, we have
$G(\varphi_U)=p^{-1}(U)$, and consequently $G(\varphi)=X$.
\end{env}

\begin{proposition}[3.5.2]
\label{II.3.5.2}
\begin{enumerate}
  \item[\rm(i)] If $\mathcal{M}$ is a quasi-coherent graded
    $\mathcal{S}$-module, there exists a canonical functorial isomorphism
    of the $\mathcal{O}_{X'}$-module $(\mathcal{M}_{[\varphi]})^{\sim}$ onto
    the $\mathcal{O}_{X'}$-module
    $\Phi_*(\widetilde{\mathcal{M}}|G(\varphi))$.
  \item[\rm(ii)] If $\mathcal{M}'$ is a quasi-coherent graded
    $\mathcal{S}'$-module, there exists a canonical functorial homomorphism
    $\nu$ from the $(\mathcal{O}_X|G(\varphi))$-module
    $\Phi^*(\widetilde{\mathcal{M}'})$ to the
    $(\mathcal{O}_X|G(\varphi))$-module
    $(\mathcal{M}'\otimes_{\mathcal{S}'}\mathcal{S})^{\sim}|G(\varphi)$. If
    $\mathcal{S}'$ is generated by $\mathcal{S}'_1$, $\nu$ is an
    isomorphism.
\end{enumerate}
\end{proposition}

\begin{proof}
The homomorphisms in question have indeed already been defined when $Y$ is
affine (\hyperref[II.2.8.7]{2.8.7} and
\hyperref[II.2.8.8]{2.8.8}), and in the general case it suffices to verify
that they are compatible with the operations of restriction from an affine
open subset of $Y$ to a smaller open subset, which follows immediately from
the commutativity of (\hyperref[II.3.5.1.1]{3.5.1.1}).
\end{proof}

In particular, for every $n\in\mathbf{Z}$, we have a canonical homomorphism
\[
\label{II.3.5.2.1}
  \Phi^*(\mathcal{O}_{X'}(n))\longrightarrow
  \mathcal{O}_X(n)|G(\varphi).
\tag{3.5.2.1}
\]

\begin{proposition}[3.5.3]
\label{II.3.5.3}
Let $Y$, $Y'$ be two preschemes, $\psi:Y'\to Y$ a morphism,
$\mathcal{S}$ a quasi-coherent graded $\mathcal{O}_Y$-algebra, and set
$\mathcal{S}'=\psi^*(\mathcal{S})$. Then the $Y'$-scheme
$X'=\operatorname{Proj}(\mathcal{S}')$ is canonically identified with
$\operatorname{Proj}(\mathcal{S})\times_Y Y'$. Moreover, if $\mathcal{M}$ is
a quasi-coherent graded $\mathcal{S}$-module, the $\mathcal{O}_{X'}$-module
$(\psi^*(\mathcal{M}))^{\sim}$ is identified with
$\widetilde{\mathcal{M}}\otimes_Y\mathcal{O}_{Y'}$.
\end{proposition}

\begin{proof}
Note in the first place that $\psi^*(\mathcal{S})$ and
$\psi^*(\mathcal{M})$ are quasi-coherent $\mathcal{O}_{Y'}$-modules,
as are their homogeneous components
(\hyperref[0.5.1.4]{0, 5.1.4}). Let $U$ be an affine open subset of $Y$,
$U'\subset\psi^{-1}(U)$ an affine open subset of $Y'$, and $A$, $A'$ the rings
of $U$ and $U'$ respectively; we then have
$\mathcal{S}|U=\widetilde{S}$, where $S$ is a graded $A$-algebra, and
$\mathcal{S}'|U'$ is identified with $(S\otimes_A A')^{\sim}$
(\hyperref[I.1.6.5]{I, 1.6.5}); the first assertion then follows from
(\hyperref[II.2.8.10]{2.8.10}) and from
(\hyperref[I.3.2.6.2]{I, 3.2.6.2}), because we verify at once that the
projection
$\operatorname{Proj}(\mathcal{S}'|U')\to
\operatorname{Proj}(\mathcal{S}|U)$ defined by the preceding identification
is compatible with the operations of restriction on $U$ and $U'$ and
therefore does define a morphism
$\operatorname{Proj}(\mathcal{S}')\to\operatorname{Proj}(\mathcal{S})$.
Now let
\[
  q:\operatorname{Proj}(\mathcal{S})\to Y,
  \qquad q':\operatorname{Proj}(\mathcal{S}')\to Y'
\]
be the structural morphisms; $q'^{-1}(U')$ is then identified with
$q^{-1}(U)\times_U U'$, and the two sheaves
$(\psi^*(\mathcal{M}))^{\sim}|q'^{-1}(U')$ and
$(\widetilde{\mathcal{M}}\otimes_Y\mathcal{O}_{Y'})|q'^{-1}(U')$
are then both canonically identified with $(M\otimes_A A')^{\sim}$, where
we have set $M=\Gamma(U,\mathcal{M})$, by virtue of
(\hyperref[II.2.8.10]{2.8.10}) and of
(\hyperref[I.1.6.5]{I, 1.6.5}); whence the second assertion, because we
again verify immediately the compatibility of the preceding identifications
with the operations of restriction.
\end{proof}

\begin{corollary}[3.5.4]
\label{II.3.5.4}
With the notation of (\hyperref[II.3.5.3]{3.5.3}),
$\mathcal{O}_{X'}(n)$ is canonically identified with
$\mathcal{O}_X(n)\otimes_Y\mathcal{O}_{Y'}$ for every $n\in\mathbf{Z}$ (with
$X=\operatorname{Proj}(\mathcal{S})$).
\end{corollary}

\begin{proof}
Indeed, with the notation of (\hyperref[II.3.5.3]{3.5.3}), it is clear
that $\psi^*(\mathcal{S}(n))=\mathcal{S}'(n)$ for every $n\in\mathbf{Z}$.
\end{proof}

\begin{env}[3.5.5]
\label{II.3.5.5}
Keeping the preceding notation, let us denote by $\Psi$ the canonical
projection $X'\to X$, and set $\mathcal{M}'=\psi^*(\mathcal{M})$; we
will moreover suppose that $\mathcal{S}$ is generated by $\mathcal{S}_1$

\oldpage[II]{63}

and that $X$ is of finite type over $Y$; it then follows that $\mathcal{S}'$
is generated by $\mathcal{S}'_1$ (as we see by reducing to the case where
$Y$ and $Y'$ are affine) and that $X'$ is of finite type over $Y'$
(\hyperref[I.6.3.4]{I, 6.3.4}). Let $\mathcal{F}$ be an
$\mathcal{O}_X$-module and set $\mathcal{F}'=\Psi^*(\mathcal{F})$; it
then follows from (\hyperref[II.3.5.4]{3.5.4}) and from
(\hyperref[0.4.3.3]{0, 4.3.3}) that we have
$\mathcal{F}'(n)=\Psi^*(\mathcal{F}(n))$ for every $n\in\mathbf{Z}$. We
moreover define a canonical $\psi$-homomorphism
$\theta_n:q_*(\mathcal{F}(n))\to q'_*(\mathcal{F}'(n))$ in the following
way: in view of the commutativity of the diagram
\[
  \xymatrix{
    X\ar[d]_{q}
    & X'\ar[l]_{\Psi}\ar[d]^{q'}\\
    Y
    & Y'\ar[l]^{\psi}
  }
\]
we have to define a homomorphism
$q_*(\mathcal{F}(n))\to
\psi_*(q'_*(\Psi^*(\mathcal{F}(n))))
=q_*(\Psi_*(\Psi^*(\mathcal{F}(n))))$, and it suffices to take
the homomorphism $\theta_n=q_*(\rho_n)$, where $\rho_n$ is the canonical
homomorphism
$\mathcal{F}(n)\to\Psi_*(\Psi^*(\mathcal{F}(n)))$
(\hyperref[0.4.4.3]{0, 4.4.3}). It is immediate that for every affine open
subset $U$ of $Y$ and every affine open subset $U'$ of $Y'$ such that
$U'\subset\psi^{-1}(U)$, the homomorphism $\theta_n$ gives on sections
the canonical homomorphism (\hyperref[0.3.7.2]{0, 3.7.2})
$\Gamma(q^{-1}(U),\mathcal{F}(n))\to
\Gamma(q'^{-1}(U'),\mathcal{F}'(n))$.

The commutativity of (\hyperref[II.2.8.13.2]{2.8.13.2}) then shows that if
$\mathcal{F}$ is quasi-coherent, the diagram
\[
  \xymatrix{
    \mathcal{F}\ar[r]^{\rho}
    & \mathcal{F}'\\
    (\Gamma_*(\mathcal{F}))^{\sim}
      \ar[u]^{\beta_{\mathcal{F}}}\ar[r]_{\widetilde{\theta}}
    & (\Gamma_*(\mathcal{F}'))^{\sim}
      \ar[u]_{\beta_{\mathcal{F}'}}
  }
\]
is commutative (the top horizontal arrow being the canonical
$\Psi$-morphism $\mathcal{F}\to\Psi^*(\mathcal{F})$).

Similarly, the commutativity of
(\hyperref[II.2.8.13.1]{2.8.13.1}) shows that the diagram
\[
  \xymatrix{
    \Gamma_*(\widetilde{\mathcal{M}})\ar[r]^{\theta}
    & \Gamma_*(\widetilde{\mathcal{M}'})\\
    \mathcal{M}\ar[r]_{\rho}\ar[u]^{\alpha_{\mathcal{M}}}
    & \mathcal{M}'\ar[u]_{\alpha_{\mathcal{M}'}}
  }
\]
is commutative (the bottom horizontal arrow being the canonical
$\psi$-morphism $\mathcal{M}\to\psi^*(\mathcal{M})$).
\end{env}

\begin{env}[3.5.6]
\label{II.3.5.6}
Now consider two preschemes $Y$, $Y'$, a morphism
$g:Y'\to Y$, a quasi-coherent graded $\mathcal{O}_Y$-algebra (resp.
$\mathcal{O}_{Y'}$-algebra) $\mathcal{S}$ (resp.
$\mathcal{S}'$) and a $g$-morphism of graded algebras
$u:\mathcal{S}\to\mathcal{S}'$, that is, an
$\mathcal{O}_Y$-homomorphism of graded algebras
$\mathcal{S}\to g_*(\mathcal{S}')$; we know moreover that it amounts to the
same thing to give an $\mathcal{O}_{Y'}$-homomorphism of graded algebras
$u^\sharp:g^*(\mathcal{S})\to\mathcal{S}'$. We then deduce canonically from
$u^\sharp$ a $Y'$-morphism
$W=\operatorname{Proj}(u^\sharp):G(u^\sharp)\to
\operatorname{Proj}(g^*(\mathcal{S}))$, where $G(u^\sharp)$ is an open subset
of $X'=\operatorname{Proj}(\mathcal{S}')$
(\hyperref[II.3.5.1]{3.5.1}). On the other hand,
$X''=\operatorname{Proj}(g^*(\mathcal{S}))$ is canonically identified

\oldpage[II]{64}

with $X\times_Y Y'$, where we set
$X=\operatorname{Proj}(\mathcal{S})$
(\hyperref[II.3.5.3]{3.5.3}); composing with
$\operatorname{Proj}(u^\sharp)$ the first projection
$p:X\times_Y Y'\to X$, we therefore obtain a morphism
$v:G(u^\sharp)\to X$, which we will denote by $\operatorname{Proj}(u)$, and
which is such that the diagram
\[
  \xymatrix{
    G(u^\sharp)\ar[r]^{v}\ar[d]
    & X\ar[d]\\
    Y'\ar[r]_{g}
    & Y
  }
\]
is commutative.

Moreover, for every quasi-coherent graded $\mathcal{S}$-module $\mathcal{M}$,
we have a canonical $v$-morphism
\[
\label{II.3.5.6.1}
  \nu:\widetilde{\mathcal{M}}\longrightarrow
  (g^*(\mathcal{M})\otimes_{g^*(\mathcal{S})}\mathcal{S}')^{\sim}
  |G(u^\sharp).
\tag{3.5.6.1}
\]

Indeed, $\nu^\sharp$ is obtained by composing the homomorphisms
\[
  v^*(\widetilde{\mathcal{M}})
  =W^*(p^*(\widetilde{\mathcal{M}}))
  \longrightarrow W^*((g^*(\mathcal{M}))^{\sim})
  \longrightarrow
  (g^*(\mathcal{M})\otimes_{g^*(\mathcal{S})}\mathcal{S}')^{\sim}
  |G(u^\sharp)
\]
where the first arrow comes from the isomorphism
(\hyperref[II.3.5.3]{3.5.3}) and the second is the homomorphism
(\hyperref[II.3.5.2]{3.5.2, (ii)}); when $\mathcal{S}$ is generated
by $\mathcal{S}_1$, it follows from (\hyperref[II.3.5.2]{3.5.2}) that
$\nu^\sharp$ is an isomorphism.

As a particular case of (\hyperref[II.3.5.6.1]{3.5.6.1}), we have, for
every $n\in\mathbf{Z}$, a canonical $v$-morphism
\[
\label{II.3.5.6.2}
  \nu:\mathcal{O}_X(n)\longrightarrow
  \mathcal{O}_{X'}(n)|G(u^\sharp).
\tag{3.5.6.2}
\]
\end{env}

\subsection{Closed sub-preschemes of a prescheme
$\operatorname{Proj}(\mathcal{S})$.}
\label{subsection:II.3.6}

\begin{env}[3.6.1]
\label{II.3.6.1}
Let $Y$ be a prescheme, and
$\varphi:\mathcal{S}\to\mathcal{S}'$ a homomorphism of quasi-coherent graded
$\mathcal{O}_Y$-algebras, of degree $0$. We say that
$\varphi$ is (\textbf{TN})-\emph{surjective} (resp.
(\textbf{TN})-\emph{injective}, (\textbf{TN})-\emph{bijective}) if there exists $n$
such that, for $k\geq n$, $\varphi_k:\mathcal{S}_k\to\mathcal{S}'_k$ is
surjective (resp. injective, bijective). When this is so, we reduce the study
of the corresponding morphism
$\Phi:\operatorname{Proj}(\mathcal{S}')\to\operatorname{Proj}(\mathcal{S})$
to the case where $\varphi$ is surjective (resp. injective, bijective).
This is proved as in (\hyperref[II.2.9.1]{2.9.1}) (which is its
particular case corresponding to $Y$ affine) by using
(\hyperref[II.3.1.8]{3.1.8}). Instead of saying that $\varphi$ is
(\textbf{TN})-bijective, we also say that it is then a
(\textbf{TN})-isomorphism.
\end{env}

\begin{proposition}[3.6.2]
\label{II.3.6.2}
Let $Y$ be a prescheme, $\mathcal{S}$ a quasi-coherent graded
$\mathcal{O}_Y$-algebra, and let $X=\operatorname{Proj}(\mathcal{S})$.
\begin{enumerate}
  \item[\rm(i)] If $\varphi:\mathcal{S}\to\mathcal{S}'$ is a
    (\textbf{TN})-surjective homomorphism of graded $\mathcal{O}_Y$-algebras,
    the corresponding morphism
    $\Phi=\operatorname{Proj}(\varphi)$
    (\hyperref[II.3.5.1]{3.5.1}) is defined on the whole of
    $\operatorname{Proj}(\mathcal{S}')$ and is a closed
    immersion of $\operatorname{Proj}(\mathcal{S}')$ into $X$. If $\mathcal{J}$
    is the kernel of $\varphi$, the closed sub-prescheme of $X$ associated to
    $\Phi$ is defined by the quasi-coherent sheaf of ideals
    $\widetilde{\mathcal{J}}$ of $\mathcal{O}_X$.
  \item[\rm(ii)] Suppose moreover that $\mathcal{S}_0=\mathcal{O}_Y$,
    that $\mathcal{S}$ is generated by $\mathcal{S}_1$ and that $\mathcal{S}_1$
    is of finite type. Let $X'$ be a closed sub-prescheme of
    $X=\operatorname{Proj}(\mathcal{S})$, defined by a quasi-coherent sheaf
    $\mathcal{I}$ of ideals of $\mathcal{O}_X$. Let
    $\mathcal{J}$

\oldpage[II]{65}

    be the quasi-coherent graded sheaf of ideals of $\mathcal{S}$ that is the
    inverse image of $\Gamma_*(\mathcal{I})$ under the canonical homomorphism
    $\alpha:\mathcal{S}\to\Gamma_*(\mathcal{O}_X)$
    (\hyperref[II.3.3.2]{3.3.2}), and set
    $\mathcal{S}'=\mathcal{S}/\mathcal{J}$. Then $X'$ is the sub-prescheme
    associated (\hyperref[I.4.2.1]{I, 4.2.1}) to the closed immersion
    $\operatorname{Proj}(\mathcal{S}')\to X$ corresponding to the canonical
    homomorphism $\mathcal{S}\to\mathcal{S}'$ of graded
    $\mathcal{O}_Y$-algebras.
\end{enumerate}
\end{proposition}

\begin{proof}
\begin{enumerate}
  \item[\rm(i)] We may suppose $\varphi$ surjective
    (\hyperref[II.3.6.1]{3.6.1}). Then, for every affine open subset $U$ of
    $Y$, $\Gamma(U,\mathcal{S})\to\Gamma(U,\mathcal{S}')$ is surjective
    (\hyperref[I.1.3.9]{I, 1.3.9}), hence
    (\hyperref[II.3.5.1]{3.5.1}) we have $G(\varphi)=X$. We are
    immediately reduced to proving the proposition when $Y$ is affine,
    and then it follows from (\hyperref[II.2.9.2]{2.9.2, (i)}).
  \item[\rm(ii)] We are reduced to proving that the homomorphism
    $\widetilde{\mathcal{J}}\to\mathcal{O}_X$ deduced from the canonical injection
    $\mathcal{J}\to\mathcal{S}$ is an isomorphism of
    $\widetilde{\mathcal{J}}$ onto $\mathcal{I}$; since the question is local
    on $Y$, we may suppose $Y$ affine with ring $A$, which implies
    $\mathcal{S}=\widetilde{S}$, where $S$ is a graded $A$-algebra generated
    by $S_1$, $S_1$ being of finite type over $A$. It then suffices to apply
    (\hyperref[II.2.9.2]{2.9.2, (ii)}).
\end{enumerate}
\end{proof}

\begin{corollary}[3.6.3]
\label{II.3.6.3}
Under the conditions of (\hyperref[II.3.6.2]{3.6.2, (i)}), and supposing
moreover that $\mathcal{S}$ is generated by $\mathcal{S}_1$,
$\Phi^*(\mathcal{O}_X(n))$ is canonically identified with
$\mathcal{O}_{X'}(n)$ for every $n\in\mathbf{Z}$.
\end{corollary}

\begin{proof}
We have defined such a canonical isomorphism when $Y$ is affine
(\hyperref[II.2.9.3]{2.9.3}); in the general case, it suffices to verify
that the isomorphisms thus defined for each affine open subset $U$ of $Y$ are
compatible with the passage from $U$ to an affine open subset $U'\subset U$,
which is immediate.
\end{proof}

\begin{corollary}[3.6.4]
\label{II.3.6.4}
Let $Y$ be a prescheme, $\mathcal{S}$ a quasi-coherent graded
$\mathcal{O}_Y$-algebra generated by $\mathcal{S}_1$, $\mathcal{M}$ a
quasi-coherent $\mathcal{O}_Y$-module, $u$ a surjective
$\mathcal{O}_Y$-homomorphism $\mathcal{M}\to\mathcal{S}_1$, and
$\overline{u}:\mathbf{S}_{\mathcal{O}_Y}(\mathcal{M})\to\mathcal{S}$
the homomorphism of graded $\mathcal{O}_Y$-algebras extending $u$
(\hyperref[II.1.7.4]{1.7.4}). The morphism corresponding to
$\overline{u}$ is then a closed immersion of
$\operatorname{Proj}(\mathcal{S})$ into
$\operatorname{Proj}(\mathbf{S}_{\mathcal{O}_Y}(\mathcal{M}))$.
\end{corollary}

\begin{proof}
Indeed, $\overline{u}$ is surjective by hypothesis and we apply
(\hyperref[II.3.6.1]{3.6.1, (i)}).
\end{proof}

\subsection{Morphisms from a prescheme to a homogeneous spectrum.}
\label{subsection:II.3.7}

\begin{env}[3.7.1]
\label{II.3.7.1}
Let $q:X\to Y$ be a morphism of preschemes, $\mathcal{L}$ an invertible
$\mathcal{O}_X$-module, and $\mathcal{S}$ a quasi-coherent $\mathcal{O}_Y$-algebra
graded in positive degrees; $q^*(\mathcal{S})$ is then a quasi-coherent
$\mathcal{O}_X$-algebra graded in positive degrees.
Consider the quasi-coherent graded $\mathcal{O}_X$-algebra
$\mathcal{S}'=\bigoplus_{n\geq0}\mathcal{L}^{\otimes n}$ and suppose given
an $\mathcal{O}_X$-homomorphism of graded algebras
\[
  \psi:q^*(\mathcal{S})\longrightarrow
  \mathcal{S}'=\bigoplus_{n\geq0}\mathcal{L}^{\otimes n}
\]
which, moreover, amounts to giving a $q$-morphism of graded algebras
\[
  \psi^\flat:\mathcal{S}\longrightarrow q_*(\mathcal{S}').
\]
We know that $\operatorname{Proj}(\mathcal{S}')$ is canonically identified with
$X$ (\hyperref[II.3.1.7]{3.1.7} and
\hyperref[II.3.1.8]{3.1.8, (iii)}); from $\psi$ we canonically deduce an
open set $G(\psi)$ of $X$ and a $Y$-morphism
\[
\label{II.3.7.1.1}
  r_{\mathcal{L},\psi}:G(\psi)\longrightarrow
  \operatorname{Proj}(\mathcal{S})=P
\tag{3.7.1.1}
\]

\oldpage[II]{66}

which we will call the morphism \emph{associated to $\mathcal{L}$ and to $\psi$};
recall (\hyperref[II.3.5.6]{3.5.6}) that this morphism is obtained by
composing with the first projection
$\pi:\operatorname{Proj}(q^*(\mathcal{S}))=P\times_YX\to P$ the
$Y$-morphism
\[
  \tau=\operatorname{Proj}(\psi):G(\psi)\longrightarrow
  \operatorname{Proj}(q^*(\mathcal{S})).
\]
\end{env}

\begin{env}[3.7.2]
\label{II.3.7.2}
Let us make $r=r_{\mathcal{L},\psi}$ explicit when $Y=\operatorname{Spec}(A)$ is
affine, and consequently $\mathcal{S}=\widetilde{S}$, where $S$ is an
$A$-algebra graded in positive degrees. Suppose first that
$X=\operatorname{Spec}(B)$ is also affine and that we have
$\mathcal{L}=\widetilde{L}$, where $L$ is a free
$B$-module of rank $1$. We then have
$q^*(\mathcal{S})=(S\otimes_AB)^{\sim}$
(\hyperref[I.1.6.5]{I, 1.6.5}); if $c$ is a generator of $L$,
$\psi_n:q^*(\mathcal{S}_n)\to\mathcal{L}^{\otimes n}$ corresponds to a
homomorphism
$w_n:s\otimes b\mapsto bv_n(s)c^{\otimes n}$ from $S_n\otimes_AB$ to
$L^{\otimes n}$, where $v_n:S_n\to B$ is a homomorphism of $A$-modules, the
$v_n$ constituting an algebra homomorphism $S\to B$. Let $f\in S_d$
($d>0$) and set $g=v_d(f)$; we have
$\pi^{-1}(D_+(f))=D_+(f\otimes1)$ by virtue of
(\hyperref[II.2.8.10]{2.8.10}) and of the identification of $D_+(f)$ with
$\operatorname{Spec}(S_{(f)})$ (\hyperref[II.2.3.6]{2.3.6}); on the other
hand, formula (\hyperref[II.2.8.1.1]{2.8.1.1}) shows (taking into account
the canonical identification of $X$ and $\operatorname{Proj}(\mathcal{S}')$)
that
\[
  \tau^{-1}(D_+(f\otimes1))=D(g)
\]
whence
\[
\label{II.3.7.2.1}
  r^{-1}(D_+(f))=D(g).
\tag{3.7.2.1}
\]

Moreover, the morphism $\tau=\operatorname{Proj}(\psi)$, restricted to $D(g)$,
corresponds to the homomorphism that maps
$(s\otimes1)/(f\otimes1)^n$ (for $s\in S_{nd}$) to
$v_{nd}(s)/g^n$ (\hyperref[II.2.8.1]{2.8.1}), and the projection $\pi$,
restricted to $D_+(f\otimes1)$, corresponds to the homomorphism
$s/f^n\mapsto(s\otimes1)/(f\otimes1)^n$; we conclude that $r$, restricted to
$D(g)$, corresponds to the homomorphism of $A$-algebras
$\omega:S_{(f)}\to B_g$ such that $\omega(s/f^n)=v_{nd}(s)/g^n$ for
$s\in S_{nd}$ ($n>0$). Passing to the case where $X$ is arbitrary ($Y$ still
being affine), we therefore obtain, taking into account
(\hyperref[II.2.8.1]{2.8.1}):
\end{env}

\begin{proposition}[3.7.3]
\label{II.3.7.3}
If $Y=\operatorname{Spec}(A)$ is affine and $\mathcal{S}=\widetilde{S}$, where
$S$ is a graded $A$-algebra, then for every
$f\in S_d=\Gamma(Y,\mathcal{S}_d)$ we have
\[
\label{II.3.7.3.1}
  r_{\mathcal{L},\psi}^{-1}(D_+(f))=X_{\psi^\flat(f)}
  \qquad\bigl(\text{where }\psi^\flat(f)\in
  \Gamma(X,\mathcal{L}^{\otimes d})\bigr)
\tag{3.7.3.1}
\]
and the restriction $X_{\psi^\flat(f)}\to D_+(f)=\operatorname{Spec}(S_{(f)})$
of $r_{\mathcal{L},\psi}$ corresponds
(\hyperref[I.2.2.4]{I, 2.2.4}) to the algebra homomorphism
\[
\label{II.3.7.3.2}
  \psi_{(f)}^\flat:S_{(f)}\longrightarrow
  \Gamma(X_{\psi^\flat(f)},\mathcal{O}_X)
\tag{3.7.3.2}
\]
such that, for $s\in S_{nd}=\Gamma(Y,\mathcal{S}_{nd})$,
\[
\label{II.3.7.3.3}
  \psi_{(f)}^\flat(s/f^n)
  =(\psi^\flat(s)|X_{\psi^\flat(f)})
   (\psi^\flat(f)|X_{\psi^\flat(f)})^{-n}.
\tag{3.7.3.3}
\]
\end{proposition}

We will say that $r_{\mathcal{L},\psi}$ is \emph{everywhere defined} if we have
$G(\psi)=X$. For this to be so, it is evidently necessary and sufficient that
$G(\psi)\cap q^{-1}(U)=q^{-1}(U)$ for every affine open subset $U\subset Y$;
in other words, the question is \emph{local on $Y$}. If $Y$ is affine,
$G(\psi)$ is the union of the $r^{-1}(D_+(f))$ for $f$ homogeneous in $S_+$
(\hyperref[II.2.8.1]{2.8.1}); by virtue of
(\hyperref[II.3.7.3.1]{3.7.3.1}), the $X_{\psi^\flat(f)}$ must therefore
form a \emph{covering} of $X$; in other words:

\begin{corollary}[3.7.4]
\label{II.3.7.4}
Under the hypotheses of (\hyperref[II.3.7.3]{3.7.3}), for
$r_{\mathcal{L},\psi}$ to be everywhere defined, it is necessary and sufficient
that for every $x\in X$ there exist an integer $n>0$ and a section $s$ of
$\mathcal{S}_n$ over $Y$ such that, if we set
$t=\psi^\flat(s)\in\Gamma(X,\mathcal{L}^{\otimes n})$, we have $t(x)\neq0$.
\end{corollary}

\oldpage[II]{67}

Note that this condition is always satisfied if $\psi$ is
(\textbf{TN})-surjective.

Similarly, the question of whether $r_{\mathcal{L},\psi}$ is \emph{dominant}
is local on $Y$, and we have:

\begin{corollary}[3.7.5]
\label{II.3.7.5}
Under the hypotheses of (\hyperref[II.3.7.3]{3.7.3}), for
$r_{\mathcal{L},\psi}$ to be dominant, it is necessary and sufficient that, for
every integer $n>0$, every section $s\in S_n$ such that
$\psi^\flat(s)\in\Gamma(X,\mathcal{L}^{\otimes n})$ is locally
nilpotent be itself nilpotent.
\end{corollary}

\begin{proof}
Indeed, we must express that
$r_{\mathcal{L},\psi}^{-1}(D_+(s))$ is empty only if $D_+(s)$ is empty,
and the corollary therefore follows from
(\hyperref[II.3.7.3.1]{3.7.3.1}) and from
(\hyperref[II.2.3.7]{2.3.7}).
\end{proof}

\begin{proposition}[3.7.6]
\label{II.3.7.6}
Let $q:X\to Y$ be a morphism, $\mathcal{L}$ an invertible
$\mathcal{O}_X$-module, $\mathcal{S}$, $\mathcal{S}'$ two quasi-coherent graded
$\mathcal{O}_Y$-algebras,
$u:\mathcal{S}'\to\mathcal{S}$ a homomorphism of graded algebras,
$\psi:q^*(\mathcal{S})\to\bigoplus_{n\geq0}\mathcal{L}^{\otimes n}$ a
homomorphism of graded algebras, and
$\psi'=\psi\circ q^*(u)$ the composite homomorphism. If
$r_{\mathcal{L},\psi'}$ is everywhere defined, so is
$r_{\mathcal{L},\psi}$; if $u$ is (\textbf{TN})-surjective and if
$r_{\mathcal{L},\psi'}$ is dominant, so is
$r_{\mathcal{L},\psi}$; conversely, if $u$ is (\textbf{TN})-injective and if
$r_{\mathcal{L},\psi}$ is dominant, $r_{\mathcal{L},\psi'}$ is dominant.
\end{proposition}

\begin{proof}
Indeed, we have $G(\psi')\subset G(\psi)$
(\hyperref[II.2.8.4]{2.8.4}), whence the first assertion; if $u$ is
(\textbf{TN})-surjective,
$\operatorname{Proj}(u):\operatorname{Proj}(\mathcal{S})\to
\operatorname{Proj}(\mathcal{S}')$ is everywhere defined and is a closed
immersion; since $r_{\mathcal{L},\psi'}$ is the composite of
$\operatorname{Proj}(u)$ and of the restriction of $r_{\mathcal{L},\psi}$ to
$G(\psi')$, we conclude that if $r_{\mathcal{L},\psi'}$ is dominant, so is
$r_{\mathcal{L},\psi}$. Finally, if $u$ is
(\textbf{TN})-injective, we know that $\operatorname{Proj}(u)$ is a dominant
morphism from $G(u)$ to $\operatorname{Proj}(\mathcal{S}')$
(\hyperref[II.2.8.3]{2.8.3}); since $G(\psi')$ is the inverse image of
$G(u)$ under $r_{\mathcal{L},\psi}$, we see that if $r_{\mathcal{L},\psi}$ is
dominant, so is $r_{\mathcal{L},\psi'}$.
\end{proof}

\begin{proposition}[3.7.7]
\label{II.3.7.7}
Let $Y$ be a quasi-compact prescheme, $q:X\to Y$ a quasi-compact morphism,
$\mathcal{L}$ an invertible $\mathcal{O}_X$-module, and $\mathcal{S}$ a
quasi-coherent graded $\mathcal{O}_Y$-algebra which is the filtered inductive
limit of an inductive system $(\mathcal{S}^\lambda)$ of quasi-coherent
$\mathcal{O}_Y$-algebras. Let
$\varphi_\lambda:\mathcal{S}^\lambda\to\mathcal{S}$ be the canonical
homomorphism,
$\psi:q^*(\mathcal{S})\to\bigoplus_{n\geq0}\mathcal{L}^{\otimes n}$ a
homomorphism of graded algebras, and set
$\psi_\lambda=\psi\circ q^*(\varphi_\lambda)$. For
$r_{\mathcal{L},\psi}$ to be everywhere defined, it is necessary and sufficient
that there exist a $\lambda$ such that $r_{\mathcal{L},\psi_\lambda}$ is
everywhere defined; $r_{\mathcal{L},\psi_\mu}$ is then everywhere defined for
$\mu\geq\lambda$.
\end{proposition}

\begin{proof}
The condition is sufficient by virtue of
(\hyperref[II.3.7.6]{3.7.6}). Conversely, suppose
$r_{\mathcal{L},\psi}$ everywhere defined; we may reduce to the case where $Y$
is affine, because if for every affine open subset $U\subset Y$ there exists
$\lambda(U)$ such that the restriction of $r_{\mathcal{L},\psi_{\lambda(U)}}$
to $q^{-1}(U)$ is everywhere defined, it will suffice ($Y$ being
quasi-compact) to cover $Y$ by finitely many affine open subsets $U_i$, and
to take $\lambda\geq\lambda(U_i)$ for all the indices $i$, by virtue of
(\hyperref[II.3.7.6]{3.7.6}). If $Y$ is affine, the hypothesis implies that
for every $x\in X$ there is a section $s^{(x)}$ of some $S_n$ such that, if
we set $t^{(x)}=\psi^\flat(s^{(x)})$, we have $t^{(x)}(x)\neq0$
($t^{(x)}$ being considered as a section of $\mathcal{L}^{\otimes n}$
over $X$), which implies $t^{(x)}(z)\neq0$ for every $z$ in a
neighbourhood $V(x)$ of $x$. Let us cover $X$ by finitely many $V(x_i)$ and
let $s^{(i)}$ be the corresponding sections of $S$; there then exists a
$\lambda$ such that the $s^{(i)}$ are all of the form
$\varphi_\lambda(s_\lambda^{(i)})$, with $s_\lambda^{(i)}\in S^\lambda$ for
all $i$; it therefore follows from (\hyperref[II.3.7.4]{3.7.4}) that
$r_{\mathcal{L},\psi_\lambda}$ is everywhere defined. The last assertion is a
trivial consequence of (\hyperref[II.3.7.6]{3.7.6}).
\end{proof}

\begin{corollary}[3.7.8]
\label{II.3.7.8}
Under the hypotheses of (\hyperref[II.3.7.7]{3.7.7}), if the
$r_{\mathcal{L},\psi_\lambda}$ are dominant, so is
$r_{\mathcal{L},\psi}$; the converse is true when the
$\varphi_\lambda$ are injective.
\end{corollary}

\oldpage[II]{68}

\begin{proof}
The second assertion is a particular case of
(\hyperref[II.3.7.6]{3.7.6}); on the other hand, to show that
$r_{\mathcal{L},\psi}$ is dominant, we may restrict to the case where $Y$ is
affine; if $s\in S$ is such that $\psi^\flat(s)$ is locally
nilpotent, then, since we can write $s=\varphi_\lambda(s_\lambda)$ for at
least one $\lambda$, we conclude from the hypothesis and from
(\hyperref[II.3.7.5]{3.7.5}) that $s_\lambda$ is nilpotent, hence so is $s$,
and the criterion of (\hyperref[II.3.7.5]{3.7.5}) applies.
\end{proof}

\begin{namedtheorem}[3.7.9]{Remarks}
\label{II.3.7.9}
\begin{enumerate}
  \item[(i)] With the notation of (\hyperref[II.3.7.1]{3.7.1}) and
    taking into account (\hyperref[II.3.2.10]{3.2.10}), we have, for every
    $n\in\mathbf{Z}$, a canonical homomorphism
    \[
    \label{II.3.7.9.1}
      \theta:r_{\mathcal{L},\psi}^*(\mathcal{O}_P(n))\longrightarrow
      \mathcal{L}^{\otimes n}
    \tag{3.7.9.1}
    \]
    defined in a general way in
    (\hyperref[II.3.5.6.2]{3.5.6.2}). We see at once that, under the
    conditions of (\hyperref[II.3.7.3]{3.7.3}), the restriction of this
    homomorphism to $X_{\psi^\flat(f)}$ is given explicitly as assigning
    to $s/f^k$ ($s\in S_{n+kd}$) the element
    $(\psi^\flat(s)|X_{\psi^\flat(f)})
    (\psi^\flat(f)|X_{\psi^\flat(f)})^{-k}$.
  \item[(ii)] Let $\mathcal{F}$ be a quasi-coherent $\mathcal{O}_X$-module, and
    suppose that $q$ is quasi-compact and separated, so that for every
    $n\geq0$, $q_*(\mathcal{F}\otimes\mathcal{L}^{\otimes n})$ is a
    quasi-coherent $\mathcal{O}_Y$-module
    (\hyperref[I.9.2.2]{I, 9.2.2}). Let
    $\mathcal{M}'=\bigoplus_{n\geq0}\mathcal{F}\otimes
    \mathcal{L}^{\otimes n}$, which is a quasi-coherent graded
    $\mathcal{S}'$-module, and consider its image
    $\mathcal{M}=q_*(\mathcal{M}')=
    \bigoplus_{n\geq0}q_*(\mathcal{F}\otimes\mathcal{L}^{\otimes n})$
    (which is a quasi-coherent $\mathcal{S}$-module, by means of
    the homomorphism $\psi^\flat$). We will see that there is a
    canonical homomorphism of $\mathcal{O}_X$-modules
    \[
    \label{II.3.7.9.2}
      \xi:r_{\mathcal{L},\psi}^*(\widetilde{\mathcal{M}})
      \longrightarrow\mathcal{F}|G(\psi).
    \tag{3.7.9.2}
    \]

    Indeed, we have defined (\hyperref[II.3.5.6.1]{3.5.6.1}) a
    canonical homomorphism
    \[
    \label{II.3.7.9.3}
      r_{\mathcal{L},\psi}^*(\widetilde{\mathcal{M}})
      \longrightarrow
      (q^*(\mathcal{M})\otimes_{q^*(\mathcal{S})}\mathcal{S}')^{\sim}
      |G(\psi),
    \tag{3.7.9.3}
    \]
    where the right-hand side is considered as a quasi-coherent sheaf on
    $\operatorname{Proj}(\mathcal{S}')$. On the other hand, we have a
    canonical homomorphism
    \[
    \label{II.3.7.9.4}
      q^*(q_*(\mathcal{M}'))\otimes_{q^*(\mathcal{S})}\mathcal{S}'
      \longrightarrow\mathcal{M}'
    \tag{3.7.9.4}
    \]
    for every quasi-coherent graded $\mathcal{S}'$-module $\mathcal{M}'$:
    for every open subset $U$ of $X$, every section $t'$ of
    $q^*(q_*(\mathcal{M}'_h))$ over $U$ and every section $b'$ of
    $\mathcal{S}'_k$ over $U$, we assign to $t'\otimes b'$ the
    section $b'\sigma(t')$ of $\mathcal{M}'_{h+k}$, where $\sigma(t')$ is the
    section of $\mathcal{M}'_h$ over $U$ which corresponds canonically
    (\hyperref[0.4.4.3]{0, 4.4.3}) to $t'$. We deduce from this a
    canonical homomorphism
    \[
    \label{II.3.7.9.5}
      (q^*(q_*(\mathcal{M}'))\otimes_{q^*(\mathcal{S})}\mathcal{S}')^{\sim}
      |G(\psi)\longrightarrow\widetilde{\mathcal{M}'}|G(\psi)
    \tag{3.7.9.5}
    \]
    and since, finally, $\widetilde{\mathcal{M}'}$ is canonically identified
    with $\mathcal{F}$
    (\hyperref[II.3.2.9]{3.2.9, (i)}), we do obtain the desired canonical
    homomorphism.

    Under the conditions of (\hyperref[II.3.7.3]{3.7.3}), the restriction
    of this homomorphism to $X_{\psi^\flat(f)}$ is given explicitly as
    follows: given a section $t_{nd}$ of
    $\mathcal{F}\otimes\mathcal{L}^{\otimes nd}$ over $X$, if
    $t'_{nd}$ is the section $t_{nd}$ considered as a section of
    $q_*(\mathcal{F}\otimes\mathcal{L}^{\otimes nd})$ over $Y$, to
    the element $t'_{nd}/f^n$ there corresponds the section
    $(t_{nd}|X_{\psi^\flat(f)})
    (\psi^\flat(f)|X_{\psi^\flat(f)})^{-n}$ of $\mathcal{F}$ over
    $X_{\psi^\flat(f)}$.
\end{enumerate}
\end{namedtheorem}

\oldpage[II]{69}


\subsection{Criteria for immersion into a homogeneous spectrum.}
\label{subsection:II.3.8}

\begin{env}[3.8.1]
\label{II.3.8.1}
With the notation of (\hyperref[II.3.7.1]{3.7.1}), the question of whether
$r_{\mathcal{L},\psi}$ is an immersion (resp. an open immersion, a
closed immersion) is evidently \emph{local on $Y$}.
\end{env}

\begin{proposition}[3.8.2]
\label{II.3.8.2}
Under the hypotheses of (\hyperref[II.3.7.3]{3.7.3}), for
$r_{\mathcal{L},\psi}$ to be everywhere defined and an immersion, it is
necessary and sufficient that there exist a family of sections
$s_\alpha\in S_{n_\alpha}$ ($n_\alpha>0$) such that, if we set
$f_\alpha=\psi^\flat(s_\alpha)$, the following conditions are satisfied:
\begin{enumerate}
  \item[(i)] The $X_{f_\alpha}$ form a covering of $X$.
  \item[(ii)] The $X_{f_\alpha}$ are affine open subsets.
  \item[(iii)] For every $\alpha$ and every
    $t\in\Gamma(X_{f_\alpha},\mathcal{O}_X)$, there exist an integer $m>0$ and
    an $s\in S_{mn_\alpha}$ such that
    $t=(\psi^\flat(s)|X_{f_\alpha})(f_\alpha|X_{f_\alpha})^{-m}$.
\end{enumerate}

For $r_{\mathcal{L},\psi}$ to be everywhere defined and an open
immersion, it is necessary and sufficient that there exist a family
$(s_\alpha)$ satisfying conditions (i), (ii), (iii) and:
\begin{enumerate}
  \item[(iv)] For every $n>0$ and every $s\in S_{nn_\alpha}$ such that
    $\psi^\flat(s)|X_{f_\alpha}=0$, there exists an integer $k>0$ such that
    $s_\alpha^ks=0$.
\end{enumerate}

For $r_{\mathcal{L},\psi}$ to be everywhere defined and a closed
immersion, it is necessary and sufficient that there exist a family
$(s_\alpha)$ satisfying (i), (ii), (iii) and:
\begin{enumerate}
  \item[(v)] The $D_+(s_\alpha)$ form a covering of
    $P=\operatorname{Proj}(S)$.
\end{enumerate}
\end{proposition}

\begin{proof}
Indeed, for $r$ to be an immersion (resp. a closed immersion), it is
necessary and sufficient that there exist a covering of $r(G(\psi))$ (resp. of
$P$) by sets $D_+(s_\alpha)$ such that, if we set
$V_\alpha=r^{-1}(D_+(s_\alpha))$, the restriction of $r$ to $U_\alpha$ is
a closed immersion of $V_\alpha$ into $D_+(s_\alpha)$
(\hyperref[I.4.2.4]{I, 4.2.4}). Condition (i) expresses both that
$r$ is everywhere defined and that the $D_+(s_\alpha)$ cover $r(X)$,
by (\hyperref[II.3.7.3.1]{3.7.3.1}); since $D_+(s_\alpha)$ is
affine, conditions (ii) and (iii) express that the restriction of $r$ to
$X_{f_\alpha}$ is a closed immersion into $D_+(s_\alpha)$, by
(\hyperref[I.4.2.3]{I, 4.2.3}); finally, since (iii) and (iv) express that
$\psi_{(s_\alpha)}^\flat$ is bijective (notation of
(\hyperref[II.3.7.3.2]{3.7.3.2})), (ii), (iii) and (iv) express that the
restriction of $r$ to $X_{f_\alpha}$ is an isomorphism onto
$D_+(s_\alpha)$ for every $\alpha$, hence (i), (ii), (iii) and (iv) express
that $r$ is an open immersion.
\end{proof}

\begin{corollary}[3.8.3]
\label{II.3.8.3}
Under the hypotheses of (\hyperref[II.3.7.6]{3.7.6}), if
$r_{\mathcal{L},\psi'}$ is everywhere defined and is an immersion, the same is
true of $r_{\mathcal{L},\psi}$. If we suppose moreover that $u$ is
(\textbf{TN})-surjective and if $r_{\mathcal{L},\psi'}$ is an open (resp.
closed) immersion, the same is true of $r_{\mathcal{L},\psi}$.
\end{corollary}

\begin{proof}
Indeed, by (\hyperref[II.3.8.2]{3.8.2}), there is a family of
$s'_\alpha\in S'_{n_\alpha}$ such that, if we set
$f_\alpha=\psi'^\flat(s'_\alpha)$, conditions (i), (ii), (iii) are
satisfied. Now, if we set $s_\alpha=u(s'_\alpha)$, we also have
$f_\alpha=\psi^\flat(s_\alpha)$, and if we have
$t=(\psi'^\flat(s')|X_{f_\alpha})(f_\alpha|X_{f_\alpha})^{-m}$, we also have
$t=(\psi^\flat(s)|X_{f_\alpha})(f_\alpha|X_{f_\alpha})^{-m}$ by setting
$s=u(s')$, whence the first assertion. The second follows immediately from
the fact that $\operatorname{Proj}(u)$ is a closed immersion.
\end{proof}

\begin{proposition}[3.8.4]
\label{II.3.8.4}
Suppose that the hypotheses of
(\hyperref[II.3.7.7]{3.7.7}) are satisfied, and moreover that $q:X\to Y$ is a
morphism of finite type. Then, for $r_{\mathcal{L},\psi}$ to be everywhere
defined and an immersion, it is necessary and sufficient that there exist
$\lambda$ such that $r_{\mathcal{L},\psi_\lambda}$ is everywhere defined and
is an immersion, and in this case $r_{\mathcal{L},\psi_\mu}$ is everywhere
defined and is an immersion for $\mu\geq\lambda$.
\end{proposition}

\oldpage[II]{70}

\begin{proof}
Taking into account (\hyperref[II.3.8.3]{3.8.3}), it suffices to prove that
if $r_{\mathcal{L},\psi}$ is everywhere defined and is an immersion, the same
is true of $r_{\mathcal{L},\psi_\lambda}$ for at least one $\lambda$. By the
same reasoning as in (\hyperref[II.3.7.7]{3.7.7}), using the
quasi-compactness of $Y$, we first reduce to the case where $Y$ is affine.
Since $X$ is quasi-compact, (\hyperref[II.3.8.2]{3.8.2}) shows the
existence of a \emph{finite} family $(s_i)$ of elements of $S$
($s_i\in S_{n_i}$) having properties (i), (ii) and (iii). The morphism
$X_{f_i}\to Y$ (where $f_i=\psi^\flat(s_i)$) is of finite type: indeed, it
is a morphism of affine schemes, hence it is quasi-compact
(\hyperref[I.6.6.1]{I, 6.6.1}), and it is locally of finite type since $q$
is of finite type (\hyperref[I.6.3.2]{I, 6.3.2}), and the conclusion follows
from (\hyperref[I.6.6.3]{I, 6.6.3}). The ring $B_i$ of $X_{f_i}$ is
consequently an $A$-algebra of finite type (\hyperref[I.6.3.3]{I, 6.3.3});
let $(t_{ij})$ be a family of generators of this algebra. By hypothesis
there are elements $s'_{ij}\in S_{m_{ij}n_i}$ such that
\[
  t_{ij}=(\psi^\flat(s'_{ij})|X_{f_i})
  (\psi^\flat(s_i)|X_{f_i})^{-m_{ij}}.
\]
By hypothesis there are a $\lambda$ and elements
$s_{i\lambda}\in S_{n_i}^\lambda$,
$s'_{ij\lambda}\in S_{m_{ij}n_i}^\lambda$ whose images under
$\varphi_\lambda$ are $s_i$ and $s'_{ij}$ respectively; it is clear that
$r_{\mathcal{L},\psi_\lambda}$ satisfies conditions (i), (ii) and (iii) of
(\hyperref[II.3.8.2]{3.8.2}).
\end{proof}

\begin{proposition}[3.8.5]
\label{II.3.8.5}
Let $Y$ be a quasi-compact scheme, or a prescheme whose underlying space
is Noetherian, $q:X\to Y$ a morphism of finite type, $\mathcal{L}$ an
invertible $\mathcal{O}_X$-module, $\mathcal{S}$ a quasi-coherent graded
$\mathcal{O}_Y$-algebra, and
$\psi:q^*(\mathcal{S})\to\bigoplus_{n\geq0}\mathcal{L}^{\otimes n}$ a
homomorphism of graded algebras. For $r_{\mathcal{L},\psi}$ to be
everywhere defined and an immersion, it is necessary and sufficient that
there exist an integer $n>0$ and a quasi-coherent sub-$\mathcal{O}_Y$-module of
finite type $\mathcal{E}$ of $\mathcal{S}_n$ such that:
\begin{enumerate}
  \item[a)] the homomorphism
    $\psi_n\circ q^*(j_n):q^*(\mathcal{E})\to\mathcal{L}^{\otimes n}$
    (where $j_n:\mathcal{E}\to\mathcal{S}$ is the canonical injection) is
    surjective;
  \item[b)] if we denote by $\mathcal{S}'$ the (graded)
    sub-$\mathcal{O}_Y$-algebra of $\mathcal{S}$ generated by $\mathcal{E}$, and
    by $\psi'$ the homomorphism $\psi\circ q^*(j')$ ($j'$ being the injection
    of $\mathcal{S}'$ into $\mathcal{S}$), then $r_{\mathcal{L},\psi'}$ is
    everywhere defined and is an immersion.
\end{enumerate}

When this is so, every quasi-coherent sub-$\mathcal{O}_Y$-module
$\mathcal{E}'$ of $\mathcal{S}_n$ containing $\mathcal{E}$ has the same
property, and so does the image of $\bigotimes^k\mathcal{E}$ in
$\mathcal{S}_{kn}$, for every $k>0$.
\end{proposition}

\begin{proof}
The sufficiency of the condition and the last two assertions are
particular cases of (\hyperref[II.3.8.3]{3.8.3}), taking into account the
canonical isomorphism between $\operatorname{Proj}(\mathcal{S})$ and
$\operatorname{Proj}(\mathcal{S}^{(k)})$
(\hyperref[II.3.1.8]{3.1.8}).

Let $(U_i)$ be a finite open covering of $Y$ by affine open subsets, and
set $A_i=A(U_i)$. Since $q^{-1}(U_i)$ is compact, the hypothesis that
$r_{\mathcal{L},\psi}$ is an everywhere defined immersion and
(\hyperref[II.3.8.2]{3.8.2}) imply the existence of a finite family
$(s_{ij})$ of elements of
$S^{(i)}=\Gamma(U_i,\mathcal{S})$ ($s_{ij}\in S_{n_{ij}}^{(i)}$) having
properties (i), (ii), (iii). We see as in the proof of
(\hyperref[II.3.8.4]{3.8.4}) that the morphism $X_{f_{ij}}\to U_i$ (where
$f_{ij}=\psi^\flat(s_{ij})$) is of finite type, and the ring $B_{ij}$ of
$X_{f_{ij}}$ is consequently an $A_i$-algebra of finite type, having a system
of generators of the form
$(\psi^\flat(t_{ijk})|X_{f_{ij}})(f_{ij}|X_{f_{ij}})^{-m_{ijk}}$, where
$t_{ijk}\in S_{m_{ijk}n_{ij}}^{(i)}$. Let $n$ be a common multiple of the
$m_{ijk}n_{ij}$; replacing (for each $(i,j,k)$) $s_{ij}$ by a
power $s_{ij}^\rho$ such that $\rho m_{ijk}n_{ij}=n$, and multiplying
$t_{ijk}$ by $s_{ij}^{\rho-m_{ijk}}$, we may suppose that for each $i$
all the $s_{ij}$ and $t_{ijk}$ belong to $S_n^{(i)}$ and that
$m_{ijk}=1$. Let $E_i$ be the sub-$A_i$-module of $S_n^{(i)}$ generated by these
elements (for fixed $i$). There exists a coherent sub-$\mathcal{O}_Y$-module of
finite type

\oldpage[II]{71}

$\mathcal{E}_i$ of $\mathcal{S}_n$ such that
$\mathcal{E}_i|U_i=(E_i)^{\sim}$
(\hyperref[I.9.4.7]{I, 9.4.7}). It is clear that the
sub-$\mathcal{O}_Y$-module $\mathcal{E}$ of $\mathcal{S}_n$ that is the sum of
the $\mathcal{E}_i$ answers the question (each section $f_{ij}$ being such
that, for every $x\in X_{f_{ij}}$, there is an affine neighbourhood
$V\subset X_{f_{ij}}$ of $x$ such that $f|V$ is a basis of
$\Gamma(V,\mathcal{L}^{\otimes n})$).
\end{proof}

